\documentclass[11pt,a4paper]{article}

\usepackage[utf8]{inputenc}
\usepackage[T1]{fontenc}
\usepackage[italian]{babel}
\usepackage{amsmath,amssymb,amsthm,mathtools}
\usepackage{booktabs}
\usepackage[margin=2.5cm]{geometry}
\usepackage{xcolor}
\usepackage{tikz}
\usetikzlibrary{arrows.meta,positioning,calc,decorations.pathreplacing,babel}
\usepackage[hidelinks]{hyperref}

\setcounter{secnumdepth}{2}

% ---------------------------------------------------------------------------
% Ambienti
% ---------------------------------------------------------------------------
\theoremstyle{definition}
\newtheorem{definizione}{Definizione}[section]
\newtheorem{ipotesi}[definizione]{Ipotesi}
\theoremstyle{plain}
\newtheorem{teorema}[definizione]{Teorema}
\newtheorem{lemma}[definizione]{Lemma}
\newtheorem{proposizione}[definizione]{Proposizione}
\theoremstyle{remark}
\newtheorem{osservazione}[definizione]{Osservazione}

% Riquadro per le idee chiave
\newcommand{\ideachiave}[1]{%
  \par\medskip\noindent
  \fcolorbox{blue!50!black}{blue!4}{\begin{minipage}{0.96\textwidth}
  \smallskip\textbf{Idea chiave.}\ #1\smallskip\end{minipage}}\par\medskip}

% ---------------------------------------------------------------------------
% Simboli
% ---------------------------------------------------------------------------
\newcommand{\R}{\mathbb{R}}
\newcommand{\Zn}{\mathbb{Z}_{\ge 0}}
\newcommand{\Xc}{\mathcal{X}}
\newcommand{\Uc}{\mathcal{U}}
\newcommand{\Xf}{\mathcal{X}_f}
\newcommand{\SSet}{\mathcal{SS}}
\newcommand{\CS}{\mathcal{CS}}
\newcommand{\Mc}{\mathcal{M}}
\DeclareMathOperator{\Conv}{Conv}
\newcommand{\JLt}{J^{\mathrm{LMPC},j}_{t\to t+N}}
\newcommand{\JL}{J^{\mathrm{LMPC},j}_{0\to N}}
\newcommand{\Jstar}{J^{*}_{0\to\infty}}
\newcommand{\eqp}[1]{\textup{[P-#1]}}   % equazione #1 del paper
\newcommand{\libro}[1]{\textup{[BBM, #1]}} % rimando al libro di Borrelli, Bemporad, Morari
\newcommand{\vett}[1]{\mathbf{#1}}

\title{Learning Model Predictive Control per compiti iterativi\\[4pt]
\large Guida alla lettura di U.~Rosolia e F.~Borrelli,\\
\emph{Learning Model Predictive Control for Iterative Tasks:\\
A Computationally Efficient Approach for Linear System},\\
IFAC PapersOnLine 50-1 (2017), pp.~3142--3147}
\author{}
\date{}

\begin{document}
\maketitle
\tableofcontents
\newpage

% ===========================================================================
\section{Di che cosa parla il paper}\label{sec:intro}
% ===========================================================================

Molti sistemi di controllo eseguono \emph{lo stesso compito più volte}: un
robot ripete lo stesso movimento di presa, un'auto da corsa percorre lo stesso
circuito giro dopo giro, un reattore chimico ripete lo stesso lotto di
produzione. Ogni esecuzione del compito si chiama \emph{iterazione} (o
\emph{prova}). Se le iterazioni partono tutte dallo stesso stato iniziale e
hanno lo stesso obiettivo, è naturale chiedersi se il controllore possa
\emph{imparare} dalle esecuzioni passate e fare meglio a ogni nuova
esecuzione.

Il paper propone un controllore di questo tipo, chiamato \emph{Learning Model
Predictive Control} (LMPC). È un MPC del tutto ordinario nella struttura (a
ogni istante risolve un problema di ottimo su un orizzonte finito e applica il
primo ingresso), con una differenza: \textbf{l'insieme terminale e il costo
terminale non sono progettati a priori, ma costruiti con i dati delle
iterazioni precedenti}. Più precisamente:
\begin{itemize}
  \item l'insieme terminale è l'\emph{inviluppo convesso} di tutti gli stati
        visitati nelle iterazioni passate (il \emph{convex safe set});
  \item il costo terminale è un'interpolazione convessa dei costi
        \emph{effettivamente pagati} a partire da quegli stati (la
        \emph{funzione baricentrica}).
\end{itemize}
Con queste due scelte il paper dimostra quattro proprietà:
\begin{enumerate}
  \item[(i)] \emph{stabilità asintotica}: a ogni iterazione il sistema in
        anello chiuso converge all'equilibrio desiderato $x_F$;
  \item[(ii)] \emph{fattibilità ricorsiva}: se il problema di ottimo è
        risolvibile all'istante iniziale, lo resta per sempre, e i vincoli su
        stato e ingresso non vengono mai violati;
  \item[(iii)] \emph{miglioramento delle prestazioni}: il costo pagato in
        un'iterazione non è mai maggiore di quello dell'iterazione precedente;
  \item[(iv)] \emph{ottimalità globale}: se le traiettorie convergono a una
        traiettoria limite al crescere del numero di iterazioni, quella
        traiettoria è la soluzione ottima del problema a orizzonte
        \emph{infinito}.
\end{enumerate}
Le proprietà (i) e (ii) sono i requisiti standard di ogni MPC ben progettato;
(iii) e (iv) sono quelle tipiche dell'apprendimento.

Il LMPC era già stato proposto dagli stessi autori per sistemi non lineari
(Rosolia e Borrelli, 2016). In quella versione l'insieme terminale è un
insieme \emph{discreto} di punti, e imporre che lo stato finale predetto cada
su uno di quei punti richiede variabili intere: il problema da risolvere a ogni
istante è un \emph{mixed-integer program}, molto costoso. Il contributo di
questo paper è mostrare che, \textbf{quando il sistema è lineare}, si può
sostituire l'insieme di punti con il suo inviluppo convesso senza perdere
nessuna garanzia: il problema diventa convesso (un QP, se il costo è
quadratico) e si risolve in tempi di diversi ordini di grandezza inferiori.

\ideachiave{Un MPC è stabile e ricorsivamente fattibile se ha un insieme
terminale \emph{controllo-invariante} e un costo terminale che si comporta da
\emph{funzione di Lyapunov} su quell'insieme. Il LMPC non progetta questi due
oggetti: li ricava dai dati. Il motivo per cui funziona è una proprietà dei
sistemi lineari con vincoli e costi convessi: \emph{ogni combinazione convessa
di traiettorie realizzate è ancora una traiettoria ammissibile, e il suo costo
non supera la combinazione convessa dei costi}. Tutto il paper discende da
questa osservazione.}

\paragraph{Come è organizzato questo documento.}
La sezione~\ref{sec:richiami} richiama le basi teoriche che servono: insiemi e
funzioni convesse, problemi di ottimo convessi, controllo ottimo a orizzonte
infinito, invarianza e stabilità secondo Lyapunov, le condizioni classiche di
stabilità dell'MPC, il controllo iterativo. Le sezioni dalla~\ref{sec:problema}
alla~\ref{sec:convergenza} seguono il paper nell'ordine in cui è scritto,
spiegando ogni definizione e ricostruendo le dimostrazioni passo per passo. La
sezione~\ref{sec:esempio} commenta l'esempio numerico del paper in modo
qualitativo. La sezione~\ref{sec:libro} mette il paper in relazione con il
libro di Borrelli, Bemporad e Morari, \emph{Predictive Control for Linear and
Hybrid Systems} (Cambridge University Press, 2017), di cui Borrelli, secondo
autore del paper, è il primo autore: mostra da quali risultati del libro
discendono i ragionamenti del paper e che cosa il paper aggiunge. La
sezione~\ref{sec:limiti} discute ipotesi e limiti, la
sezione~\ref{sec:riepilogo} riassume.

\paragraph{Convenzioni.}
Le equazioni del paper sono citate come \eqp{n}: \eqp{18} è l'equazione (18)
del paper. I rimandi al libro sono scritti come \libro{Teorema 12.2}, con la
numerazione del libro. Le equazioni di questo documento hanno la numerazione
normale. Il
paper usa la lettera $k$ con due significati diversi (indice di iterazione
nelle sommatorie e istante di predizione nel problema di MPC) e la lettera $t$
sia per l'istante corrente sia come indice muto nelle sommatorie sui dati. Qui
i significati sono separati come nella tabella~\ref{tab:notazione}.

\begin{table}[h]
\centering
\caption{Notazione usata nel documento.}\label{tab:notazione}
\begin{tabular}{ll}
\toprule
Simbolo & Significato \\
\midrule
$x\in\R^n$, $u\in\R^m$ & stato e ingresso del sistema \\
$\Xc$, $\Uc$ & insiemi (convessi) dei valori ammessi per stato e ingresso \\
$x_S$, $x_F$ & stato iniziale di ogni iterazione, equilibrio finale \\
$h(x,u)$ & costo di stadio \\
$j$ & iterazione corrente \\
$i$ & iterazione generica (passata), $i\le j$ \\
$t$ & istante corrente all'interno di un'iterazione \\
$s$ & istante a cui è stato registrato un dato, $x^i_s$, $u^i_s$ \\
$k$ & istante predetto: $x_{k|t}$ è la predizione di $x_k$ fatta all'istante $t$ \\
$l$ & indice muto nelle somme dei costi \\
$N$ & orizzonte di predizione \\
$x^i_s$, $u^i_s$ & stato e ingresso \emph{realizzati} all'istante $s$ dell'iterazione $i$ \\
$J^i_{s\to\infty}(x^i_s)$ & costo realizzato da $x^i_s$ in poi (\emph{cost-to-go}) \\
$\Jstar(x)$ & costo ottimo a orizzonte infinito a partire da $x$ \\
$\SSet^j$, $\CS^j$ & safe set campionato e safe set convesso dopo l'iterazione $j$ \\
$\lambda^i_s$ & peso del dato $x^i_s$ in una combinazione convessa \\
$P^j(x)$ & funzione baricentrica (costo terminale) \\
$\JLt(x^j_t)$ & costo ottimo del problema LMPC all'istante $t$ dell'iterazione $j$ \\
\bottomrule
\end{tabular}
\end{table}

% ===========================================================================
\section{Richiami teorici}\label{sec:richiami}
% ===========================================================================

\subsection{Sistemi lineari vincolati}\label{sec:sistemi}

Il paper considera sistemi lineari a tempo discreto e tempo invarianti
\begin{equation}\label{eq:sistema}
  x_{t+1} = A x_t + B u_t, \qquad x_t\in\R^n,\; u_t\in\R^m,
\end{equation}
soggetti ai vincoli
\begin{equation}\label{eq:vincoli}
  x_t\in\Xc,\qquad u_t\in\Uc\qquad \forall t\in\Zn,
\end{equation}
con $\Xc$ e $\Uc$ insiemi \emph{convessi}. Nella pratica sono quasi sempre
poliedri, cioè insiemi descritti da disuguaglianze lineari, ad esempio limiti
superiori e inferiori su ogni componente ($\underline{x}\le x\le\overline{x}$,
$\underline{u}\le u\le\overline{u}$).

Un punto $x_F$ è un \emph{equilibrio del sistema non forzato} se
\begin{equation}\label{eq:equilibrio}
  x_F = A x_F,
\end{equation}
cioè se, applicando $u=0$, lo stato resta in $x_F$. Il caso più comune è
$x_F=0$. Il paper chiede che $x_F$ sia anche \emph{ammissibile}, cioè
$x_F\in\Xc$ e $0\in\Uc$: il sistema può restare in $x_F$ per sempre senza
violare i vincoli.

\subsection{Insiemi convessi, inviluppo convesso, funzioni convesse}\label{sec:convessita}

\begin{definizione}[Combinazione convessa]
Dati i punti $z_1,\dots,z_M\in\R^n$, una loro \emph{combinazione convessa} è
un punto
\[
  z = \sum_{r=1}^{M}\lambda_r z_r,\qquad \lambda_r\ge 0,\qquad
  \sum_{r=1}^{M}\lambda_r = 1.
\]
I pesi $\lambda_r$ sono anche detti \emph{coordinate baricentriche} di $z$
rispetto ai punti $z_r$.
\end{definizione}

Geometricamente, la combinazione convessa di due punti è un punto del segmento
che li unisce; quella di tre punti è un punto del triangolo che ha i tre punti
per vertici; e così via.

\begin{definizione}[Insieme convesso, inviluppo convesso]
Un insieme $C\subseteq\R^n$ è \emph{convesso} se contiene il segmento che
unisce due qualsiasi dei suoi punti. L'\emph{inviluppo convesso}
$\Conv(S)$ di un insieme $S$ è il più piccolo insieme convesso che contiene
$S$; coincide con l'insieme di tutte le combinazioni convesse (finite) di
punti di $S$. Se $S$ è un insieme finito di punti, $\Conv(S)$ è un
\emph{politopo}: un poliedro limitato i cui vertici sono alcuni dei punti di
$S$.
\end{definizione}

Un politopo si può descrivere in due modi equivalenti: con i suoi vertici
(\emph{rappresentazione V}: ``tutte le combinazioni convesse di questi
punti'') o con le disuguaglianze delle sue facce (\emph{rappresentazione H}:
$\{x: Hx\le g\}$). Passare da una all'altra è costoso in dimensione alta: il
libro \libro{Sez.~4.4.2} osserva che il calcolo dell'inviluppo convesso in
forma H ha un costo esponenziale nel numero di facce. Il
paper userà sempre la rappresentazione V, che il libro scrive in forma
compatta \libro{eq.~(4.8)} come
\[
  \Conv(V) = \{V\lambda : \lambda\ge 0,\ \vett{1}^\top\lambda = 1\},
\]
con $V$ matrice che ha per colonne i punti: non calcola mai le facce
dell'inviluppo convesso, ma tratta i pesi $\lambda$ come variabili di
ottimizzazione.

\begin{definizione}[Funzione convessa]
Una funzione $f:C\to\R$, con $C$ convesso, è \emph{convessa} se
\[
  f(\theta z_1 + (1-\theta) z_2) \le \theta f(z_1) + (1-\theta) f(z_2)
  \qquad \forall z_1,z_2\in C,\;\theta\in[0,1].
\]
Il costo di stadio $h(x,u)$ è \emph{congiuntamente convesso} se è convesso
come funzione della coppia $(x,u)$, cioè del vettore $(x,u)\in\R^{n+m}$.
\end{definizione}

L'esempio tipico è il costo quadratico $h(x,u)=x^\top Q x + u^\top R u$ con
$Q\succeq 0$ e $R\succeq 0$, congiuntamente convesso perché è una forma
quadratica semidefinita positiva del vettore $(x,u)$.

\begin{lemma}[Disuguaglianza di Jensen]\label{lem:jensen}
Se $f$ è convessa, per ogni combinazione convessa vale
\[
  f\Big(\sum_{r=1}^M \lambda_r z_r\Big) \le \sum_{r=1}^M \lambda_r f(z_r).
\]
\end{lemma}
\begin{proof}
Per $M=2$ è la definizione di convessità. Se vale per $M-1$ punti e
$\lambda_M<1$, si scrive
$\sum_{r=1}^{M}\lambda_r z_r = (1-\lambda_M)\,w + \lambda_M z_M$ con
$w=\sum_{r<M}\frac{\lambda_r}{1-\lambda_M}z_r$, che è una combinazione
convessa di $M-1$ punti. Applicando la definizione e poi l'ipotesi induttiva
a $w$ si ottiene la tesi. (Se $\lambda_M=1$ la tesi è banale.)
\end{proof}

In parole: \emph{la funzione valutata nella media pesata non supera la media
pesata dei valori}. Nel paper la disuguaglianza di Jensen è il passaggio
centrale della dimostrazione di stabilità (equazione~\eqp{31}).

\subsection{Il lemma che regge tutto il paper}\label{sec:lemma-traiettorie}

Il risultato seguente non compare nel paper come enunciato separato, ma è
usato in tutte le dimostrazioni. Conviene isolarlo subito.

\begin{lemma}[Combinazioni convesse di traiettorie]\label{lem:traiettorie}
Siano $(x^{(r)}_s, u^{(r)}_s)_{s\ge 0}$, per $r=1,\dots,M$, traiettorie del
sistema~\eqref{eq:sistema}: $x^{(r)}_{s+1} = A x^{(r)}_s + B u^{(r)}_s$, con
$x^{(r)}_s\in\Xc$ e $u^{(r)}_s\in\Uc$ per ogni $s$. Siano $\lambda_r\ge 0$ con
$\sum_r\lambda_r=1$. Allora le successioni
\[
  \tilde x_s = \sum_{r=1}^M\lambda_r x^{(r)}_s, \qquad
  \tilde u_s = \sum_{r=1}^M\lambda_r u^{(r)}_s
\]
\begin{enumerate}
  \item sono ancora una traiettoria del sistema: $\tilde x_{s+1}=A\tilde x_s+B\tilde u_s$;
  \item rispettano i vincoli: $\tilde x_s\in\Xc$, $\tilde u_s\in\Uc$;
  \item hanno un costo non maggiore della combinazione dei costi, se $h$ è
        congiuntamente convesso:
        $h(\tilde x_s,\tilde u_s)\le\sum_r\lambda_r h(x^{(r)}_s,u^{(r)}_s)$.
\end{enumerate}
\end{lemma}
\begin{proof}
(1) Per linearità,
$A\tilde x_s + B\tilde u_s = \sum_r\lambda_r(Ax^{(r)}_s+Bu^{(r)}_s)
=\sum_r\lambda_r x^{(r)}_{s+1}=\tilde x_{s+1}$.
(2) $\tilde x_s$ è combinazione convessa di punti di $\Xc$, che è convesso;
lo stesso per $\tilde u_s$ e $\Uc$.
(3) È la disuguaglianza di Jensen applicata a $h$ nei punti
$(x^{(r)}_s,u^{(r)}_s)$.
\end{proof}

\begin{osservazione}
Il punto (1) è il \emph{principio di sovrapposizione} dei sistemi lineari,
ristretto a pesi che sommano a uno. Per un sistema non lineari
$x_{s+1}=f(x_s,u_s)$ è falso: la media di due traiettorie in generale non è
una traiettoria. È esattamente questo il motivo per cui il LMPC non lineare
del 2016 non può usare l'inviluppo convesso e deve ricorrere a variabili
intere.
\end{osservazione}

\subsection{Problemi di ottimo convessi: LP, QP, MIP}\label{sec:ottimizzazione}

Un problema di ottimo $\min_z f(z)$ con $z\in Z$ è \emph{convesso} se $f$ è
convessa e $Z$ è convesso. La proprietà fondamentale è che \emph{ogni minimo
locale è anche globale}: non esistono ``valli'' secondarie in cui un
algoritmo possa rimanere intrappolato. Due casi particolari interessano qui:
\begin{itemize}
  \item \emph{programmazione lineare} (LP): costo lineare $c^\top z$, vincoli
        lineari $Gz\le g$, $Ez=e$;
  \item \emph{programmazione quadratica} (QP): costo
        $\tfrac12 z^\top H z + f^\top z$ con $H\succeq 0$, vincoli lineari.
\end{itemize}
Entrambi si risolvono in tempo polinomiale con algoritmi maturi (simplesso,
punti interni, insiemi attivi: \texttt{quadprog} in MATLAB è uno di questi).
Il libro li tratta in \libro{Cap.~2} e i metodi numerici in \libro{Cap.~3}.

Se invece alcune variabili devono essere \emph{intere} (tipicamente binarie,
$\delta\in\{0,1\}$), il problema è un \emph{mixed-integer program} (MIP,
oppure MIQP se il costo è quadratico). L'insieme ammissibile non è più
convesso: è un'unione di pezzi, uno per ogni combinazione delle variabili
binarie. Gli algoritmi esatti (\emph{branch and bound}) esplorano in modo
intelligente queste combinazioni, ma nel caso peggiore il loro numero cresce
in modo esponenziale con il numero di variabili binarie: i problemi con
variabili 0-1 sono in generale NP-completi \libro{Sez.~2.4 e~17.4.1}. Nel LMPC del 2016
c'è una variabile binaria per ogni stato memorizzato, e gli stati memorizzati
aumentano a ogni iterazione: da qui il costo computazionale che questo paper
elimina.

\subsection{Controllo ottimo a orizzonte infinito}\label{sec:ottimo-infinito}

Il problema che il paper vuole risolvere è
\begin{equation}\label{eq:ottimo-infinito}
\begin{aligned}
  \Jstar(x_S) = \min_{u_0,u_1,\dots}\ & \sum_{k=0}^{\infty} h(x_k,u_k)\\
  \text{s.t. }\ & x_{k+1} = Ax_k + Bu_k,\quad x_0 = x_S,\\
  & x_k\in\Xc,\ u_k\in\Uc\quad \forall k\ge 0.
\end{aligned}
\end{equation}
La funzione $x\mapsto\Jstar(x)$, che associa a ogni stato iniziale il costo
ottimo, si chiama \emph{funzione valore} (o \emph{costo ottimo da percorrere},
\emph{optimal cost-to-go}). Essa soddisfa il \emph{principio di ottimalità di
Bellman}:
\begin{equation}\label{eq:bellman}
  \Jstar(x) = \min_{u\in\Uc,\ Ax+Bu\in\mathcal{D}}
  \Big\{ h(x,u) + \Jstar(Ax+Bu)\Big\},
\end{equation}
dove $\mathcal{D}$ è l'insieme degli stati da cui il problema è risolvibile.
In parole: il costo ottimo da $x$ è il costo del primo passo più il costo
ottimo dallo stato in cui si arriva. Se si conoscesse $\Jstar$, basterebbe
risolvere a ogni istante il problema a un passo~\eqref{eq:bellman} per ottenere
il controllo ottimo a orizzonte infinito \libro{Sez.~7.4, eq.~(7.21)--(7.22)}.

Senza vincoli e con costo quadratico, $\Jstar(x)=x^\top P_\infty x$ con
$P_\infty$ soluzione dell'equazione algebrica di Riccati (LQR). \emph{Con} i
vincoli (\emph{constrained} LQR, CLQR) non esiste una formula chiusa: la
funzione valore è quadratica a tratti su una partizione poliedrica dello
spazio di stato, e calcolarla esattamente è possibile solo per problemi
piccoli \libro{Sez.~11.3.4--11.3.5}. L'esempio del paper (sezione~\ref{sec:esempio})
usa proprio questa soluzione esatta come termine di paragone.

\subsection{Invarianza e stabilità secondo Lyapunov}\label{sec:lyapunov}

\begin{definizione}[Insieme controllo-invariante, \libro{Def.~10.9}]\label{def:invariante}
Un insieme $C\subseteq\Xc$ è \emph{controllo-invariante} per il
sistema~\eqref{eq:sistema} con i vincoli~\eqref{eq:vincoli} se
\[
  \forall x\in C\quad \exists u\in\Uc:\quad Ax+Bu\in C.
\]
\end{definizione}
Se lo stato è in un insieme controllo-invariante, esiste sempre un ingresso
ammissibile che lo mantiene nell'insieme, e quindi lo si può mantenere
nell'insieme (e nei vincoli) \emph{per sempre}. Per questo gli insiemi
controllo-invarianti sono detti anche ``sicuri''.

Si consideri ora un sistema in anello chiuso $x_{t+1}=f(x_t)$ con un
equilibrio $x_F=f(x_F)$.

\begin{definizione}[Stabilità asintotica]
L'equilibrio $x_F$ è \emph{stabile} se per ogni $\varepsilon>0$ esiste
$\delta>0$ tale che $\|x_0-x_F\|<\delta$ implica $\|x_t-x_F\|<\varepsilon$
per ogni $t$; è \emph{asintoticamente stabile} se è stabile e inoltre
$x_t\to x_F$ per $t\to\infty$ (per stati iniziali in un opportuno insieme,
la \emph{regione di attrazione}).
\end{definizione}

\begin{teorema}[Lyapunov, tempo discreto, \libro{Teorema 7.2}]\label{teo:lyapunov}
Sia $\Omega$ un insieme positivamente invariante per $x_{t+1}=f(x_t)$ con
$x_F\in\Omega$. Se esiste una funzione $V:\Omega\to\R$ continua in $x_F$ con
\[
  V(x_F)=0,\qquad V(x)>0\ \ \forall x\ne x_F,\qquad
  V(f(x)) - V(x) \le -\alpha(x) < 0\ \ \forall x\ne x_F,
\]
con $\alpha$ continua e definita positiva, allora $x_F$ è asintoticamente
stabile con regione di attrazione $\Omega$.
\end{teorema}

L'idea è intuitiva: $V$ è una sorta di ``energia'' che è nulla solo
all'equilibrio e diminuisce strettamente a ogni passo; non potendo diminuire
per sempre di quantità non trascurabili, deve tendere a zero, e con essa la
distanza da $x_F$. Nell'MPC il ruolo di $V$ è svolto dal costo ottimo del
problema di predizione, e $\alpha(x)$ è il costo di stadio $h$.

\subsection{Il controllo predittivo e le sue garanzie classiche}\label{sec:mpc}

L'MPC rende trattabile il problema~\eqref{eq:ottimo-infinito} troncandolo a un
orizzonte finito $N$. All'istante $t$, misurato lo stato $x_t$, risolve
\begin{equation}\label{eq:mpc}
\begin{aligned}
  J^{\mathrm{MPC}}_N(x_t) = \min_{u_{t|t},\dots,u_{t+N-1|t}}\ &
  \sum_{k=t}^{t+N-1} h(x_{k|t},u_{k|t}) + V_f(x_{t+N|t})\\
  \text{s.t. }\ & x_{k+1|t} = Ax_{k|t} + Bu_{k|t},\quad x_{t|t}=x_t,\\
  & x_{k|t}\in\Xc,\ u_{k|t}\in\Uc,\quad k=t,\dots,t+N-1,\\
  & x_{t+N|t}\in\Xf,
\end{aligned}
\end{equation}
applica il primo ingresso $u_t = u^*_{t|t}$ e all'istante successivo ripete
tutto con il nuovo stato misurato (\emph{orizzonte recedente}, figura~\ref{fig:recedente}).
Gli oggetti $\Xf$ (\emph{insieme terminale}) e $V_f$ (\emph{costo terminale})
rappresentano ``ciò che succede dopo l'orizzonte''.

\begin{figure}[h]
\centering
\begin{tikzpicture}[>=Stealth, font=\small]
  \draw[->] (0,0) -- (10.6,0) node[below left]{tempo};
  \draw[->] (0,0) -- (0,3.4) node[left]{stato};
  \fill[green!15] (7.65,0) rectangle (8.35,0.9);
  \node[green!40!black] at (8,1.15) {$\Xf$};
  \draw[thick] plot[smooth] coordinates {(0.4,2.9) (1.4,2.7) (2.4,2.4) (3.4,2.1)};
  \foreach \p in {(0.4,2.9),(1.4,2.7),(2.4,2.4),(3.4,2.1)} \fill \p circle (1.6pt);
  \draw[dashed, blue, thick] (3.4,2.1) -- (4.55,1.7) -- (5.7,1.3) -- (6.85,0.95) -- (8,0.6);
  \foreach \p in {(4.55,1.7),(5.7,1.3),(6.85,0.95),(8,0.6)} \fill[blue] \p circle (1.6pt);
  \draw[dotted, red, thick] (4.55,1.7) -- (5.7,1.25) -- (6.85,0.85) -- (8,0.55) -- (9.15,0.35);
  \draw[dotted] (3.4,0) -- (3.4,3.1);
  \foreach \xx/\lab in {3.4/t, 4.55/t{+}1, 8/t{+}N}
     \draw (\xx,0.06) -- (\xx,-0.06) node[below]{$\lab$};
  \draw[decorate, decoration={brace, mirror, amplitude=5pt}] (3.4,-0.6) -- (8,-0.6)
     node[midway, below=6pt]{orizzonte di predizione $N$};
  \node[orange!80!black, font=\footnotesize, above] at (4.05,2.25) {applicato};
  \begin{scope}[xshift=8.9cm, yshift=2.9cm, font=\footnotesize]
    \draw[thick] (0,0) -- (0.5,0); \node[right] at (0.5,0) {stati misurati};
    \draw[dashed, blue, thick] (0,-0.4) -- (0.5,-0.4); \node[right] at (0.5,-0.4) {predizione in $t$};
    \draw[dotted, red, thick] (0,-0.8) -- (0.5,-0.8); \node[right] at (0.5,-0.8) {predizione in $t{+}1$};
  \end{scope}
\end{tikzpicture}
\caption{Orizzonte recedente. All'istante $t$ si predice l'evoluzione per $N$
passi, imponendo che lo stato finale cada nell'insieme terminale $\Xf$; si
applica solo il primo ingresso. All'istante $t+1$ la finestra si sposta di un
passo e il problema viene risolto di nuovo.}
\label{fig:recedente}
\end{figure}

Troncare l'orizzonte ha un prezzo. Un MPC ``ingenuo'' (senza $\Xf$ e $V_f$
adeguati) è \emph{miope}: può portare il sistema in uno stato da cui, qualche
passo dopo, non esiste più alcun ingresso che rispetti i vincoli (perdita di
fattibilità), oppure può non convergere all'equilibrio (perdita di
stabilità). La teoria classica (Mayne et al., 2000) dà condizioni sufficienti
per evitare entrambi i problemi.

\begin{ipotesi}[Condizioni classiche sugli ingredienti terminali]\label{ip:terminali}
\leavevmode
\begin{enumerate}
  \item $h(x_F,0)=0$ e $h(x,u)>0$ per ogni $(x,u)\ne(x_F,0)$;
  \item $\Xf\subseteq\Xc$ è controllo-invariante e contiene $x_F$;
  \item su $\Xf$ il costo terminale $V_f$ è una \emph{funzione di Lyapunov di
        controllo} \libro{Remark 12.3}: per ogni $x\in\Xf$ esiste $u\in\Uc$
        tale che
        \begin{equation}\label{eq:clf}
          Ax+Bu\in\Xf\qquad\text{e}\qquad V_f(Ax+Bu) - V_f(x)\le -h(x,u).
        \end{equation}
\end{enumerate}
\end{ipotesi}

\begin{teorema}[Stabilità dell'MPC, \libro{Teorema 12.2}]\label{teo:mpc}
Sotto l'ipotesi~\ref{ip:terminali}, se il problema~\eqref{eq:mpc} ha
soluzione in $t=0$, allora ha soluzione per ogni $t\ge 0$ e lungo l'anello
chiuso
\begin{equation}\label{eq:decrescita-mpc}
  J^{\mathrm{MPC}}_N(x_{t+1}) - J^{\mathrm{MPC}}_N(x_t) \le -h(x_t,u_t).
\end{equation}
Di conseguenza (con le opportune ipotesi di continuità) $x_F$ è
asintoticamente stabile, con regione di attrazione l'insieme degli stati da
cui~\eqref{eq:mpc} ha soluzione.
\end{teorema}
\begin{proof}[Idea della dimostrazione]
Sia $(u^*_{t|t},\dots,u^*_{t+N-1|t})$ la soluzione ottima all'istante $t$ e
$x^*_{t+N|t}\in\Xf$ lo stato finale predetto. Per il punto 3
dell'ipotesi esiste $\bar u$ che mantiene $x^*_{t+N|t}$ in $\Xf$. All'istante
$t+1$ si considera la sequenza \emph{traslata}
\[
  (u^*_{t+1|t},\dots,u^*_{t+N-1|t},\ \bar u):
\]
le prime $N-1$ mosse sono quelle già pianificate, l'ultima è quella fornita
dall'ipotesi. Poiché il modello è esatto, $x_{t+1}=x^*_{t+1|t}$ e la sequenza
traslata produce gli stessi stati predetti, più un ultimo stato
$\bar x = Ax^*_{t+N|t}+B\bar u\in\Xf$: è \emph{ammissibile} (questo dà la
fattibilità ricorsiva). Il suo costo è quello ottimo all'istante $t$, meno il
primo termine $h(x_t,u_t)$, meno il vecchio costo terminale, più il nuovo
termine di stadio e il nuovo costo terminale:
\[
  J^{\mathrm{MPC}}_N(x_t) - h(x_t,u_t)
  + \underbrace{h(x^*_{t+N|t},\bar u) + V_f(\bar x) - V_f(x^*_{t+N|t})}_{\le 0\ \text{per la~\eqref{eq:clf}}}.
\]
Il costo ottimo in $t+1$ non supera quello di una soluzione ammissibile, e
si ottiene la~\eqref{eq:decrescita-mpc}. Il costo ottimo è quindi una
funzione di Lyapunov e il teorema~\ref{teo:lyapunov} conclude.
\end{proof}

\ideachiave{Il teorema~\ref{teo:mpc} ha bisogno di una sola cosa ``difficile'':
una coppia $(\Xf,V_f)$ che soddisfi la~\eqref{eq:clf}. Nella pratica la si
costruisce attorno all'equilibrio con un controllore lineare locale (ad
esempio LQR) e un insieme invariante piccolo, il che rende la regione di
attrazione piccola e obbliga a orizzonti lunghi. Il LMPC usa la stessa
dimostrazione, ma con una coppia $(\Xf,V_f)$ ricavata dai dati, grande e
sempre più vicina alla coppia ideale (dominio di $\Jstar$, $\Jstar$).}

\subsection{Compiti iterativi e Iterative Learning Control}\label{sec:ilc}

L'\emph{Iterative Learning Control} (ILC) studia sistemi che ripetono lo stesso
compito. Il tempo ha due scale: l'istante $t$ \emph{all'interno} di
un'iterazione e l'indice $j$ \emph{dell'iterazione}. Le iterazioni partono
dallo stesso stato iniziale, e alla fine di ognuna i dati registrati sono
disponibili per progettare la successiva (figura~\ref{fig:scale}).

Nella forma classica l'ILC insegue un riferimento $r_t$ noto e aggiorna
l'ingresso con l'errore dell'iterazione precedente, con leggi del tipo
\[
  u^{j+1}_t = u^j_t + L\,(r_{t+1} - y^j_{t+1}),
\]
in modo che l'errore di inseguimento si annulli al crescere di $j$. Questo
presuppone che \emph{il riferimento sia noto}. In molti compiti però il
riferimento ``migliore'' è proprio ciò che non si sa: in una gara
automobilistica non si conosce a priori la traiettoria che minimizza il tempo
sul giro. Il LMPC è un controllo iterativo \emph{senza riferimento}
(\emph{reference-free}): l'obiettivo non è inseguire una traiettoria, ma
minimizzare un costo; ciò che si impara dalle iterazioni non è un segnale di
correzione, ma una descrizione sempre migliore di \emph{dove} il sistema può
stare in sicurezza e di \emph{quanto costa} arrivare alla fine da lì.

Il paper cita anche un lavoro (Tamar et al., 2016) in cui il costo terminale
di un MPC è appreso da una rete neurale: funziona in pratica, ma senza alcuna
garanzia di stabilità, fattibilità o miglioramento. Il LMPC ottiene invece
queste garanzie perché il costo terminale non è un'approssimazione generica,
ma è costruito con costi \emph{realmente pagati} lungo traiettorie
\emph{realmente percorse}.

\begin{figure}[h]
\centering
\begin{tikzpicture}[>=Stealth, font=\small]
  \foreach \j/\y in {0/0, 1/-1.4, 2/-2.8}{
    \node[circle, fill, inner sep=1.6pt, label=left:{$x_S$}] (s\j) at (0,\y) {};
    \node[circle, fill, inner sep=1.6pt, label=above:{$x_F$}] (f\j) at (7,\y) {};
    \draw[->, thick] (s\j) .. controls (2.5,{\y+0.55}) and (4.5,{\y-0.45}) .. (f\j)
      node[pos=0.5, above=4pt]{iterazione $\j$};
  }
  \node at (3.5,-3.6) {$\vdots$};
  \draw[->] (0,1.0) -- (7,1.0) node[right]{tempo $t$ (dentro l'iterazione)};
  \draw[->] (-1.2,0.3) -- (-1.2,-3.5) node[below]{iterazione $j$};
  \node[draw, rounded corners, align=center, fill=gray!8] (mem) at (10.2,-1.4)
     {memoria dei dati:\\ stati, ingressi,\\ costi realizzati};
  \foreach \j in {0,1,2} \draw[->, dashed, gray] (f\j) -- (mem.west);
\end{tikzpicture}
\caption{Le due scale temporali di un compito iterativo. Ogni iterazione parte
da $x_S$ e deve arrivare a $x_F$; a fine iterazione la traiettoria percorsa e
i costi pagati vengono memorizzati e usati per progettare le iterazioni
successive.}
\label{fig:scale}
\end{figure}

% ===========================================================================
\section{Il problema affrontato dal paper}\label{sec:problema}
% ===========================================================================

\paragraph{Dati di un'iterazione.}
All'iterazione $j$ il sistema~\eqref{eq:sistema} viene fatto evolvere e si
registrano le successioni di ingressi e stati (\eqp{3})
\[
  \vett{u}^j = [u^j_0,\,u^j_1,\,\dots,\,u^j_t,\,\dots],\qquad
  \vett{x}^j = [x^j_0,\,x^j_1,\,\dots,\,x^j_t,\,\dots].
\]
Il pedice è il tempo, l'apice è l'iterazione: $x^j_t$ è lo stato
\emph{realizzato} (misurato, non predetto) all'istante $t$ dell'iterazione
$j$. Tutte le iterazioni partono dallo stesso stato (\eqp{4}):
\begin{equation}\label{eq:stato-iniziale}
  x^j_0 = x_S\qquad\forall j\ge 0.
\end{equation}

\paragraph{Obiettivo.}
A ogni iterazione si vorrebbe risolvere il problema a orizzonte
infinito~\eqref{eq:ottimo-infinito} (\eqp{5}) a partire da $x_S$. Sul costo di
stadio il paper fa tre ipotesi (\eqp{6}):
\begin{itemize}
  \item $h$ è continuo;
  \item $h$ è congiuntamente convesso in $(x,u)$;
  \item $h$ è definito positivo rispetto a $(x_F,0)$: $h(x_F,0)=0$ e
        $h(x,u)>0$ per ogni $(x,u)\ne(x_F,0)$.\footnote{Il paper scrive
        ``$h(x^j_t,u^j_t)\succ 0$ per ogni $x^j_t\in\R^n\setminus\{x_F\}$,
        $u^j_t\in\R^m\setminus\{0\}$''; il senso inteso, quello usato nelle
        dimostrazioni, è la positività per ogni coppia diversa da
        $(x_F,0)$.}
\end{itemize}
Lo stato finale $x_F$ è un equilibrio ammissibile del sistema non forzato,
$x_F=Ax_F$ (\eqp{7}). La positività di $h$ ha una conseguenza importante: un
costo totale finito $\sum_t h(x_t,u_t)<\infty$ è possibile solo se
$h(x_t,u_t)\to 0$, cioè se la traiettoria converge a $x_F$. In altre parole,
\emph{minimizzare il costo a orizzonte infinito significa anche portare il
sistema a $x_F$}: il compito ``arrivare a $x_F$'' è codificato nel costo.

\paragraph{Che cosa si misura a ogni iterazione.}
Il problema~\eqref{eq:ottimo-infinito} non si sa risolvere direttamente (è a
orizzonte infinito e vincolato). Il LMPC lo affronta per approssimazioni
successive: a ogni iterazione applica un controllore, registra la traiettoria
e il costo effettivamente pagato, e usa questa informazione per l'iterazione
successiva. La qualità di un'iterazione si giudica dal costo realizzato
$J^j_{0\to\infty}(x_S)=\sum_{t\ge 0}h(x^j_t,u^j_t)$, definito nella
sezione~\ref{sec:costo-terminale}. Il traguardo è che questo costo scenda
iterazione dopo iterazione fino al valore ottimo $\Jstar(x_S)$.

\begin{figure}[h]
\centering
\newcommand{\disegnaDati}{%
  \draw[gray] (-0.5,0.3) rectangle (5.7,4.1);
  \node[gray, below left] at (5.7,4.1) {$\Xc$};
  \draw[thick, black!70] (0,2.5) -- (1,3.3) -- (2,3.6) -- (3,3.3) -- (3.8,2.6)
      -- (4.4,2) -- (4.8,1.7) -- (5,1.5);
  \draw[thick, teal] (0,2.5) -- (1.2,2.2) -- (2.3,1.6) -- (3.3,1.1) -- (4.2,1.1)
      -- (4.7,1.3) -- (5,1.5);
  \foreach \p in {(1,3.3),(2,3.6),(3,3.3),(3.8,2.6),(4.4,2),(4.8,1.7)}
      \fill[black!70] \p circle (1.7pt);
  \foreach \p in {(1.2,2.2),(2.3,1.6),(3.3,1.1),(4.2,1.1),(4.7,1.3)}
      \fill[teal] \p circle (1.7pt);
  \fill (0,2.5) circle (2pt) node[left]{$x_S$};
  \fill (5,1.5) circle (2pt) node[right]{$x_F$};
}
\begin{tikzpicture}[>=Stealth, font=\small, scale=0.95]
  \begin{scope}
    \disegnaDati
    \node[black!70, font=\footnotesize] at (1.6,3.85) {iterazione 0};
    \node[teal, font=\footnotesize] at (1.4,1.3) {iterazione 1};
    \node at (2.6,-0.2) {$\SSet^1$: i punti};
  \end{scope}
  \begin{scope}[xshift=7.2cm]
    \fill[blue!12] (0,2.5) -- (1,3.3) -- (2,3.6) -- (3,3.3) -- (3.8,2.6)
      -- (4.8,1.7) -- (5,1.5) -- (4.7,1.3) -- (4.2,1.1) -- (3.3,1.1) -- cycle;
    \draw[blue!60, thin] (0,2.5) -- (1,3.3) -- (2,3.6) -- (3,3.3) -- (3.8,2.6)
      -- (4.8,1.7) -- (5,1.5) -- (4.7,1.3) -- (4.2,1.1) -- (3.3,1.1) -- cycle;
    \disegnaDati
    \node at (2.6,-0.2) {$\CS^1=\Conv(\SSet^1)$};
  \end{scope}
\end{tikzpicture}
\caption{Spazio di stato (schematico, $n=2$). A sinistra: gli stati visitati
in due iterazioni riuscite formano il safe set campionato $\SSet^1$, un
insieme finito di punti. A destra: il safe set convesso $\CS^1$ è il loro
inviluppo convesso. Il rettangolo rappresenta i vincoli di stato $\Xc$.}
\label{fig:safe-set}
\end{figure}

% ===========================================================================
\section{L'insieme terminale imparato: il safe set}\label{sec:safe-set}
% ===========================================================================

\subsection{Il safe set campionato}

Si chiama \emph{iterazione riuscita} un'iterazione la cui traiettoria
converge a $x_F$. L'insieme degli indici delle iterazioni riuscite fino alla
$j$-esima è (\eqp{9})
\begin{equation}\label{eq:Mj}
  \Mc^j = \Big\{ i\in\{0,\dots,j\} : \lim_{s\to\infty} x^i_s = x_F \Big\}.
\end{equation}
Il \emph{safe set campionato} è l'insieme di tutti gli stati visitati durante
le iterazioni riuscite (\eqp{8}):
\begin{equation}\label{eq:SS}
  \SSet^j = \bigcup_{i\in\Mc^j}\ \bigcup_{s=0}^{\infty}\{x^i_s\}.
\end{equation}

\begin{proposizione}\label{prop:SS}
Ogni punto di $\SSet^j$ è ``sicuro'': da $x^i_s\in\SSet^j$ esiste una
sequenza di ingressi ammissibili che porta lo stato a $x_F$ rispettando i
vincoli e restando in $\SSet^j$. Inoltre $\SSet^j$ è controllo-invariante.
\end{proposizione}
\begin{proof}
La sequenza è quella registrata, $u^i_s,u^i_{s+1},\dots$: è già stata
applicata con successo, rispetta i vincoli, porta in $x^i_{s+1},x^i_{s+2},\dots$,
che appartengono tutti a $\SSet^j$, e converge a $x_F$ perché $i\in\Mc^j$.
\end{proof}

Si noti che per sapere che questi punti sono sicuri non serve calcolare
nulla: lo sono \emph{perché sono già stati percorsi}. È un insieme
controllo-invariante ottenuto gratuitamente dai dati, senza le difficili
procedure di calcolo degli insiemi invarianti.

Il difetto di $\SSet^j$ è che è un \emph{insieme discreto}. Usarlo come
insieme terminale significa imporre $x_{t+N|t}\in\SSet^j$, cioè che lo stato
finale predetto coincida con uno dei punti memorizzati. Nel LMPC non lineare
del 2016 questo vincolo si scrive con una variabile binaria $\delta^i_s$ per
ogni punto:
\[
  x_{t+N|t}=\sum_{i,s}\delta^i_s\,x^i_s,\qquad \delta^i_s\in\{0,1\},\qquad
  \sum_{i,s}\delta^i_s = 1,
\]
e il problema diventa un MIP con tante variabili binarie quanti sono i punti
memorizzati.

\subsection{Il safe set convesso}

L'idea del paper è \emph{rilassare} le variabili binarie in variabili
continue $\lambda^i_s\in[0,1]$. Lo stato finale non deve più coincidere con un
punto memorizzato, ma può essere una qualunque loro combinazione convessa. Si
definisce così il \emph{safe set convesso} (\eqp{10}):
\begin{equation}\label{eq:CS}
  \CS^j = \Conv(\SSet^j)
  = \Big\{ \sum_{i\in\Mc^j}\sum_{s\ge 0}\lambda^i_s x^i_s\ :\
    \lambda^i_s\ge 0,\ \sum_{i\in\Mc^j}\sum_{s\ge 0}\lambda^i_s = 1 \Big\}.
\end{equation}
(Il paper scrive la stessa cosa numerando i punti di $\SSet^j$ come
$z_1,z_2,\dots$ con pesi $\alpha_1,\alpha_2,\dots$.) Per un sistema non lineare
questo rilassamento sarebbe pericoloso: un punto ``in mezzo'' a punti sicuri
non è necessariamente sicuro. Per un sistema lineare invece lo è.

\begin{proposizione}[$\CS^j$ è controllo-invariante]\label{prop:CS-invariante}
Se $\Xc$ e $\Uc$ sono convessi, $\CS^j\subseteq\Xc$ ed è
controllo-invariante. Più precisamente, se
$x=\sum_{i,s}\lambda^i_s x^i_s\in\CS^j$, l'ingresso
\begin{equation}\label{eq:ubar}
  \bar u = \sum_{i,s}\lambda^i_s u^i_s
\end{equation}
è ammissibile e porta in
\begin{equation}\label{eq:xbar}
  Ax + B\bar u = \sum_{i,s}\lambda^i_s x^i_{s+1}\in\CS^j.
\end{equation}
\end{proposizione}
\begin{proof}
$\CS^j\subseteq\Xc$ perché $\Xc$ è convesso e contiene tutti i punti di
$\SSet^j$. Il resto è il lemma~\ref{lem:traiettorie} applicato alle traiettorie
memorizzate: $\bar u\in\Uc$ perché combinazione convessa di ingressi
ammissibili; per linearità
$Ax+B\bar u=\sum\lambda^i_s(Ax^i_s+Bu^i_s)=\sum\lambda^i_s x^i_{s+1}$; questo è
una combinazione convessa con gli \emph{stessi pesi} dei punti successori
$x^i_{s+1}$, che appartengono anch'essi a $\SSet^j$; quindi sta in $\CS^j$.
\end{proof}

La dimostrazione dice qualcosa di più dell'invarianza: \emph{insegna come
restare nell'insieme}. Se uno stato è una media pesata di stati memorizzati, si
applica la stessa media pesata degli ingressi memorizzati, e lo stato
successivo sarà la stessa media pesata degli stati successivi. Ripetendo
all'infinito, lo stato segue la ``traiettoria media'', che converge a $x_F$
perché tutte le traiettorie componenti vi convergono.

\begin{osservazione}[Monotonia]
Un'iterazione riuscita resta riuscita: $\Mc^i\subseteq\Mc^j$ per $i\le j$.
Quindi $\SSet^i\subseteq\SSet^j$ e (\eqp{11})
\begin{equation}\label{eq:CS-monotono}
  \CS^i\subseteq\CS^j\qquad\forall i\le j.
\end{equation}
L'insieme terminale può solo crescere con le iterazioni. Ogni nuova
traiettoria che esce dall'inviluppo attuale lo allarga.
\end{osservazione}

\begin{osservazione}[L'equilibrio finale]
$x_F$ è il limite delle traiettorie memorizzate. Nella pratica le
traiettorie sono finite e si memorizza come ultimo punto $x_F$ stesso (con
ingresso $u=0$ e successore ancora $x_F$), così che $x_F\in\SSet^j$. Si
assume questo nel seguito.
\end{osservazione}

% ===========================================================================
\section{Il costo terminale imparato: la funzione baricentrica}\label{sec:costo-terminale}
% ===========================================================================

\subsection{Costi realizzati}

Il secondo ingrediente è il costo terminale. Per ogni stato memorizzato
$x^i_s$ si conosce il costo che è stato \emph{effettivamente pagato} da lì in
poi lungo l'iterazione $i$ (\eqp{12}):
\begin{equation}\label{eq:cost-to-go}
  J^i_{s\to\infty}(x^i_s) = \sum_{l=0}^{\infty} h(x^i_{s+l},u^i_{s+l}).
\end{equation}
È un numero, non una funzione da ottimizzare. Si calcola a ritroso a fine
iterazione con la ricorsione
\begin{equation}\label{eq:ricorsione-costi}
  J^i_{s\to\infty}(x^i_s) = h(x^i_s,u^i_s) + J^i_{s+1\to\infty}(x^i_{s+1}),
\end{equation}
partendo dal fondo della traiettoria, dove lo stato è $x_F$ e il costo
residuo è nullo. Il costo dell'iterazione $i$ è il cost-to-go dal suo primo
stato (\eqp{13}): $J^i_{0\to\infty}(x^i_0)=J^i_{0\to\infty}(x_S)$.

Poiché ogni traiettoria memorizzata è ammissibile, il suo costo non può
essere inferiore all'ottimo:
\begin{equation}\label{eq:costo-maggiore-ottimo}
  J^i_{s\to\infty}(x^i_s)\ \ge\ \Jstar(x^i_s).
\end{equation}
Il cost-to-go realizzato è quindi una \emph{stima per eccesso} del costo
ottimo, disponibile però solo nei punti memorizzati.

\subsection{La funzione baricentrica}

Per avere un costo terminale definito su tutto $\CS^j$ e non solo nei punti
memorizzati, il paper interpola i costi realizzati. La \emph{funzione
baricentrica} (Jones e Morari, 2010) è (\eqp{14}--\eqp{15})
\begin{equation}\label{eq:P}
  P^j(x) =
  \begin{cases}
    p^{j,*}(x) & \text{se } x\in\CS^j,\\
    +\infty    & \text{se } x\notin\CS^j,
  \end{cases}
\end{equation}
dove $p^{j,*}(x)$ è il valore ottimo del programma lineare
\begin{equation}\label{eq:LP}
\begin{aligned}
  p^{j,*}(x) = \min_{\lambda^i_s\ge 0}\ & \sum_{i\in\Mc^j}\sum_{s\ge 0}
      \lambda^i_s\, J^i_{s\to\infty}(x^i_s)\\
  \text{s.t. }\ & \sum_{i\in\Mc^j}\sum_{s\ge 0}\lambda^i_s = 1,\qquad
      \sum_{i\in\Mc^j}\sum_{s\ge 0}\lambda^i_s x^i_s = x.
\end{aligned}
\end{equation}
In parole: \emph{tra tutti i modi di scrivere $x$ come media pesata di stati
memorizzati, si sceglie quello per cui la stessa media pesata dei costi
realizzati è minima}. Il valore $+\infty$ fuori da $\CS^j$ è un modo compatto
di dire che lo stato finale predetto \emph{deve} cadere in $\CS^j$: in
qualunque problema di minimo, un costo infinito equivale a un vincolo.

\begin{osservazione}[Interpretazione geometrica]
Si considerino i punti $(x^i_s,\,J^i_{s\to\infty})\in\R^{n+1}$: ogni stato
memorizzato ``sollevato'' all'altezza del suo costo. Il grafico di $P^j$ è il
\emph{bordo inferiore dell'inviluppo convesso} di questa nuvola di punti
(figura~\ref{fig:baricentrica}). È una funzione affine a tratti: sopra ogni
``faccia'' inferiore dell'inviluppo, $P^j$ interpola linearmente i costi dei
vertici della faccia. I punti che stanno sopra questo bordo (stati
memorizzati con un costo realizzato ``cattivo'') non influenzano $P^j$.
\end{osservazione}

\begin{figure}[h]
\centering
\begin{tikzpicture}[>=Stealth, font=\small, x=2.6cm, y=1cm]
  \draw[->] (-0.15,0) -- (3.6,0) node[below]{$x$};
  \draw[->] (0,-0.2) -- (0,4.3) node[left]{costo};
  \draw[dashed, thick, red!70!black, domain=0:3.3, samples=60]
     plot (\x, {0.35*\x*\x}) node[right]{$\Jstar(x)$};
  \draw[thick, blue!70!black] (0,0) -- (0.8,0.5) -- (2.2,1.9) -- (3,3.4);
  \draw[->, thick, blue!70!black] (3,3.4) -- (3,4.2) node[right, font=\footnotesize]{$+\infty$ fuori da $\CS^j$};
  \node[blue!70!black, above left] at (0.55,0.33) {$P^j(x)$};
  \foreach \p in {(0,0),(0.8,0.5),(1.1,1.4),(1.5,1.3),(2.0,2.6),(2.2,1.9),(2.6,2.9),(3,3.4)}
     \fill \p circle (2pt);
  \node[below left] at (0,0) {$x_F$};
  \draw[decorate, decoration={brace, mirror, amplitude=5pt}] (0,-0.45) -- (3,-0.45)
     node[midway, below=6pt]{$\CS^j$};
\end{tikzpicture}
\caption{Funzione baricentrica in una dimensione (schematico). I punti neri
sono le coppie $(x^i_s, J^i_{s\to\infty})$: stato memorizzato e costo
realizzato da lì in poi. $P^j$ (linea continua) è il bordo inferiore del loro
inviluppo convesso, definito su $\CS^j$. La funzione valore ottima $\Jstar$
(tratteggiata) sta sempre sotto: $P^j$ è una stima per eccesso del costo
ottimo.}
\label{fig:baricentrica}
\end{figure}

Il paper osserva (\eqp{16}) che, indicando con $\lambda^{*,i}_s$ i pesi ottimi
del problema~\eqref{eq:LP}, si può scrivere
\begin{equation}\label{eq:P-esplicita}
  P^j(x) = \sum_{i,s}\lambda^{*,i}_s J^i_{s\to\infty}(x^i_s)
  = \sum_{i,s}\lambda^{*,i}_s\sum_{l=0}^{\infty} h(x^i_{s+l},u^i_{s+l}),
\end{equation}
cioè $P^j(x)$ è la media pesata, con i pesi ottimi, dei costi pagati lungo le
traiettorie che ``passano attorno'' a $x$.

\subsection{Proprietà della funzione baricentrica}

\begin{proposizione}\label{prop:P}
La funzione $P^j$ ha le proprietà seguenti.
\begin{enumerate}
  \item[(a)] \emph{È convessa} su $\CS^j$.
  \item[(b)] \emph{Sovrastima il costo ottimo}: $P^j(x)\ge\Jstar(x)$ per ogni
        $x\in\CS^j$.
  \item[(c)] \emph{Non supera il costo realizzato} nei punti memorizzati:
        $P^j(x^i_s)\le J^i_{s\to\infty}(x^i_s)$.
  \item[(d)] \emph{Non cresce con le iterazioni}: $P^{j}(x)\le P^{j-1}(x)$
        per ogni $x$.
  \item[(e)] $P^j(x_F)=0$ e $P^j(x)\ge 0$.
\end{enumerate}
\end{proposizione}
\begin{proof}
(a) Siano $x_1,x_2\in\CS^j$ con pesi ottimi $\lambda_1$, $\lambda_2$ e sia
$\theta\in[0,1]$. I pesi $\theta\lambda_1+(1-\theta)\lambda_2$ sono non
negativi, sommano a uno e ricostruiscono $\theta x_1+(1-\theta)x_2$: sono
ammissibili per~\eqref{eq:LP} in quel punto, con costo
$\theta P^j(x_1)+(1-\theta)P^j(x_2)$. Il minimo non può essere maggiore, quindi
$P^j(\theta x_1+(1-\theta)x_2)\le\theta P^j(x_1)+(1-\theta)P^j(x_2)$. (È un
fatto generale: il valore ottimo di un LP in funzione del termine noto dei
vincoli di uguaglianza è convesso.)

(b) Con i pesi ottimi $\lambda^*$ per $x$, si combinano le traiettorie
memorizzate che partono dai punti $x^i_s$, ciascuna ``riavviata''
all'istante $s$. Per il lemma~\ref{lem:traiettorie} la combinazione è una
traiettoria ammissibile che parte da $x$, con costo non maggiore di
$\sum\lambda^{*,i}_s J^i_{s\to\infty}=P^j(x)$. Il costo ottimo $\Jstar(x)$ non
supera il costo di nessuna traiettoria ammissibile.

(c) Nel punto $x^i_s$ basta scegliere $\lambda^i_s=1$ e tutti gli altri pesi
nulli: è ammissibile e ha costo $J^i_{s\to\infty}(x^i_s)$.

(d) Passando da $j-1$ a $j$ si aggiungono punti e quindi variabili $\lambda$
al problema~\eqref{eq:LP}; i vecchi pesi restano ammissibili (con i nuovi
pesi a zero) e il minimo può solo diminuire. Fuori da $\CS^{j-1}$ il valore
passa da $+\infty$ a un valore finito o resta $+\infty$.

(e) Tutti i costi realizzati sono non negativi, quindi $P^j\ge 0$; in $x_F$
si sceglie il peso unitario sul punto memorizzato $x_F$, che ha costo
residuo nullo.
\end{proof}

La proprietà più importante è però la seguente: $P^j$ soddisfa su $\CS^j$
esattamente la condizione~\eqref{eq:clf} richiesta al costo terminale di un
MPC. Il paper la dimostra dentro la prova del teorema~1 (\eqp{30}--\eqp{31});
qui la si isola come lemma.

\begin{lemma}[$P^j$ è una funzione di Lyapunov di controllo su $\CS^j$]\label{lem:decrescita}
Sia $x\in\CS^j$ e siano $\lambda^{*,i}_s$ i pesi ottimi di~\eqref{eq:LP} in
$x$. Con l'ingresso $\bar u=\sum_{i,s}\lambda^{*,i}_s u^i_s\in\Uc$ e lo stato
successivo $\bar x=Ax+B\bar u\in\CS^j$ vale
\begin{equation}\label{eq:decrescita-P}
  P^j(\bar x) \le P^j(x) - h(x,\bar u).
\end{equation}
\end{lemma}
\begin{proof}
Che $\bar u\in\Uc$ e $\bar x=\sum\lambda^{*,i}_s x^i_{s+1}\in\CS^j$ è la
proposizione~\ref{prop:CS-invariante}. Per la disuguaglianza, si parte dalla
definizione e si usa la ricorsione~\eqref{eq:ricorsione-costi}:
\[
\begin{aligned}
  P^j(x) &= \sum_{i,s}\lambda^{*,i}_s J^i_{s\to\infty}(x^i_s)
  = \sum_{i,s}\lambda^{*,i}_s h(x^i_s,u^i_s)
   + \sum_{i,s}\lambda^{*,i}_s J^i_{s+1\to\infty}(x^i_{s+1})\\
  &\ge h\Big(\sum_{i,s}\lambda^{*,i}_s x^i_s,\ \sum_{i,s}\lambda^{*,i}_s u^i_s\Big)
   + \sum_{i,s}\lambda^{*,i}_s J^i_{s+1\to\infty}(x^i_{s+1})
   && \text{(Jensen)}\\
  &= h(x,\bar u) + \sum_{i,s}\lambda^{*,i}_s J^i_{s+1\to\infty}(x^i_{s+1})\\
  &\ge h(x,\bar u) + P^j(\bar x).
\end{aligned}
\]
L'ultimo passaggio vale perché i pesi $\lambda^{*,i}_s$, ``spostati'' sui
punti successori $x^i_{s+1}$, sono ammissibili per il problema~\eqref{eq:LP}
nel punto $\bar x$: sono non negativi, sommano a uno e ricostruiscono
$\bar x$. Il loro costo è quindi maggiore o uguale al minimo $P^j(\bar x)$.
\end{proof}

\ideachiave{Il lemma~\ref{lem:decrescita} è il cuore tecnico del paper.
Combinando le traiettorie memorizzate con i pesi che definiscono $P^j(x)$, si
ottiene un ingresso che resta in $\CS^j$ e fa diminuire $P^j$ almeno del costo
di stadio. Quindi la coppia $(\CS^j, P^j)$ soddisfa l'ipotesi~\ref{ip:terminali}
del teorema classico dell'MPC, \emph{senza averla progettata}: è stata
raccolta dai dati.}

\subsection{Iterazioni di durata finita}

Nella pratica ogni iterazione dura un tempo finito $t_i$ (\emph{Remark~3} del
paper): le somme infinite su $s$ in~\eqref{eq:CS} e~\eqref{eq:LP} diventano
somme fino a $t_i$ (\eqp{17}), e il problema~\eqref{eq:LP} diventa un LP con un
numero finito di variabili, una per ogni stato memorizzato. Se l'ultimo
stato registrato è (numericamente) $x_F$, le proprietà precedenti continuano
a valere; il criterio usato nell'esempio per decidere quando un'iterazione è
finita è discusso nella sezione~\ref{sec:esempio}.

% ===========================================================================
\section{Il controllore LMPC}\label{sec:lmpc}
% ===========================================================================

\subsection{Il problema risolto a ogni istante}

All'istante $t$ dell'iterazione $j$, misurato lo stato $x^j_t$, il LMPC
risolve (\eqp{18})
\begin{equation}\label{eq:lmpc}
\begin{aligned}
  \JLt(x^j_t) = \min_{u_{t|t},\dots,u_{t+N-1|t}}\ &
  \sum_{k=t}^{t+N-1} h(x_{k|t},u_{k|t}) + P^{j-1}(x_{t+N|t})\\
  \text{s.t. }\ & x_{k+1|t} = Ax_{k|t}+Bu_{k|t},\quad k=t,\dots,t+N-1,\\
  & x_{t|t} = x^j_t,\\
  & x_{k|t}\in\Xc,\ u_{k|t}\in\Uc,\quad k=t,\dots,t+N-1,\\
  & x_{t+N|t}\in\CS^{j-1}.
\end{aligned}
\end{equation}
È esattamente il problema MPC~\eqref{eq:mpc} con
\[
  \Xf=\CS^{j-1},\qquad V_f = P^{j-1}:
\]
insieme e costo terminale sono quelli costruiti con le iterazioni
$0,\dots,j-1$, cioè con tutti i dati disponibili all'inizio dell'iterazione
$j$. Indicate con
\[
  \vett{u}^{*,j}_{t:t+N|t}=[u^{*,j}_{t|t},\dots,u^{*,j}_{t+N-1|t}],\qquad
  \vett{x}^{*,j}_{t:t+N|t}=[x^{*,j}_{t|t},\dots,x^{*,j}_{t+N|t}]
\]
la sequenza ottima di ingressi e la traiettoria predetta corrispondente
(\eqp{19}), al sistema si applica il primo ingresso (\eqp{20}):
\begin{equation}\label{eq:legge}
  u^j_t = u^{*,j}_{t|t}.
\end{equation}
All'istante $t+1$ si misura $x^j_{t+1}$ e si risolve di nuovo~\eqref{eq:lmpc}:
è un normale orizzonte recedente. Insieme e costo terminale restano
\emph{fissi per tutta l'iterazione} $j$; vengono aggiornati solo alla fine
dell'iterazione, quando la nuova traiettoria è completa e i suoi costi
realizzati si possono calcolare con la~\eqref{eq:ricorsione-costi}.

Un aspetto da sottolineare: il vincolo terminale riguarda solo
$x_{t+N|t}$. Gli stati intermedi della predizione possono uscire da
$\CS^{j-1}$ (purché restino in $\Xc$). È così che il LMPC \emph{esplora}: può
percorrere regioni dello spazio di stato mai visitate, a patto di rientrare,
entro $N$ passi, in una zona di cui si sa che è sicura e quanto costa
(figura~\ref{fig:predizione}). La traiettoria nuova verrà memorizzata e
allargherà l'insieme terminale per le iterazioni successive.

\begin{figure}[h]
\centering
\begin{tikzpicture}[>=Stealth, font=\small, scale=1.05]
  \draw[gray] (-0.5,-0.15) rectangle (5.7,4.1);
  \node[gray, below left] at (5.7,4.1) {$\Xc$};
  \fill[blue!12] (0,2.5) -- (1,3.3) -- (2,3.6) -- (3,3.3) -- (3.8,2.6)
      -- (4.8,1.7) -- (5,1.5) -- (4.7,1.3) -- (4.2,1.1) -- (3.3,1.1) -- cycle;
  \draw[blue!60, thin] (0,2.5) -- (1,3.3) -- (2,3.6) -- (3,3.3) -- (3.8,2.6)
      -- (4.8,1.7) -- (5,1.5) -- (4.7,1.3) -- (4.2,1.1) -- (3.3,1.1) -- cycle;
  \node[blue!60!black] at (2.4,2.85) {$\CS^{j-1}$};
  \foreach \p in {(1,3.3),(2,3.6),(3,3.3),(3.8,2.6),(4.4,2),(4.8,1.7),
                  (1.2,2.2),(2.3,1.6),(3.3,1.1),(4.2,1.1),(4.7,1.3)}
      \fill[blue!40] \p circle (1.4pt);
  \fill (0,2.5) circle (2pt) node[left]{$x_S$};
  \fill (5,1.5) circle (2pt) node[right]{$x_F$};
  \draw[thick] (0,2.5) -- (0.5,1.6) -- (1.0,1.0) -- (1.5,0.6);
  \foreach \p in {(0.5,1.6),(1.0,1.0)} \fill \p circle (1.7pt);
  \fill (1.5,0.6) circle (2.2pt) node[left]{$x^j_t$};
  \draw[thick, dashed, red!80!black, ->] (1.5,0.6) -- (2.3,0.55) -- (3.0,0.75) -- (3.6,1.2);
  \foreach \p in {(2.3,0.55),(3.0,0.75)} \fill[red!80!black] \p circle (1.7pt);
  \fill[red!80!black] (3.6,1.2) circle (2.2pt);
  \draw[<-, red!80!black] (3.65,1.25) -- (4.6,0.65) node[right, align=left, font=\footnotesize]
      {$x_{t+N|t}\in\CS^{j-1}$};
  \node[font=\footnotesize, right] at (0.65,1.5) {realizzato};
  \node[red!80!black, font=\footnotesize] at (2.6,0.18) {predetto ($N=3$)};
\end{tikzpicture}
\caption{Il problema LMPC all'istante $t$ dell'iterazione $j$ (schematico). La
traiettoria predetta (tratteggiata) può uscire dal safe set convesso
$\CS^{j-1}$, ma deve rientrarvi entro $N$ passi. Il costo dello stato finale
predetto è valutato con $P^{j-1}$.}
\label{fig:predizione}
\end{figure}

\subsection{Perché è un problema convesso}\label{sec:qp}

Nel problema~\eqref{eq:lmpc} compare $P^{j-1}$, che è a sua volta il valore
ottimo di un problema di minimo. Poiché si minimizza in entrambi i casi, i due
problemi si fondono in uno solo, in cui i pesi $\lambda$ diventano variabili
di decisione insieme agli ingressi. Si indichino con $\vett{D}^{j-1}$ la
\emph{matrice dei dati}, le cui colonne sono gli stati memorizzati $x^i_s$, e
con $\vett{c}^{j-1}$ il vettore dei corrispondenti costi realizzati
$J^i_{s\to\infty}(x^i_s)$, per $i\in\Mc^{j-1}$:
\[
  \vett{D}^{j-1} = \big[\,x^0_0\ \ x^0_1\ \cdots\ x^{j-1}_{t_{j-1}}\,\big]\in\R^{n\times M},
  \qquad
  \vett{c}^{j-1} = \big[\,J^0_{0\to\infty}\ \ J^0_{1\to\infty}\ \cdots\,\big]^\top\in\R^{M},
\]
con $M$ numero totale di stati memorizzati. Il problema~\eqref{eq:lmpc}
equivale a
\begin{equation}\label{eq:lmpc-unico}
\begin{aligned}
  \min_{u_{t|t},\dots,u_{t+N-1|t},\ \lambda\in\R^M}\ &
  \sum_{k=t}^{t+N-1} h(x_{k|t},u_{k|t}) + (\vett{c}^{j-1})^\top\lambda\\
  \text{s.t. }\ & x_{k+1|t} = Ax_{k|t}+Bu_{k|t},\quad x_{t|t} = x^j_t,\\
  & x_{k|t}\in\Xc,\ u_{k|t}\in\Uc,\\
  & x_{t+N|t} = \vett{D}^{j-1}\lambda,\qquad \lambda\ge 0,\qquad \vett{1}^\top\lambda = 1.
\end{aligned}
\end{equation}
Se $h$ è quadratico e $\Xc$, $\Uc$ sono poliedri:
\begin{itemize}
  \item il costo è quadratico negli ingressi (e negli stati, che ne sono
        funzioni affini) e \emph{lineare} nei pesi $\lambda$;
  \item tutti i vincoli sono lineari: la dinamica, i limiti, il legame
        $x_{t+N|t}=\vett{D}^{j-1}\lambda$, la positività e la somma unitaria
        dei pesi.
\end{itemize}
È quindi un QP. In forma \emph{condensata}, sostituendo la dinamica, lo stato
finale predetto è
$x_{t+N|t}=A^N x^j_t + \sum_{l=0}^{N-1}A^{N-1-l}Bu_{t+l|t}$, e il vincolo
terminale diventa un vincolo di uguaglianza lineare tra ingressi e pesi:
\[
  \big[A^{N-1}B\ \ A^{N-2}B\ \cdots\ B\big]
  \begin{bmatrix}u_{t|t}\\ \vdots\\ u_{t+N-1|t}\end{bmatrix}
  - \vett{D}^{j-1}\lambda = -A^N x^j_t .
\]
L'insieme terminale non viene mai calcolato in forma esplicita (le sue facce):
è descritto implicitamente dalle colonne di $\vett{D}^{j-1}$.

\begin{osservazione}[Remark~4 del paper]
Il paper sottolinea che il problema~\eqref{eq:lmpc} è convesso perché
l'insieme terminale $\CS^{j-1}$ è convesso e il costo terminale $P^{j-1}$ è una
funzione convessa (proposizione~\ref{prop:P}(a)), e indica questa
formulazione convessa come il contributo principale rispetto al LMPC del 2016.
(Nel testo del Remark~4 i riferimenti \eqp{18a} e \eqp{18e} sono scambiati:
il vincolo terminale è \eqp{18e}, il costo terminale compare in \eqp{18a}.)
\end{osservazione}

\begin{osservazione}[Dimensione del problema]
Il numero di variabili $\lambda$ è il numero di stati memorizzati, che cresce
a ogni iterazione. Rispetto al MIP del 2016 il vantaggio è enorme (un QP con
molte variabili continue si risolve in tempo polinomiale, un MIP con molte
variabili binarie no), ma il problema continua a crescere. Nella pratica si
usano solo sottoinsiemi dei dati (ad esempio le ultime iterazioni, o i punti
vicini allo stato corrente); il paper non tratta questo aspetto.
\end{osservazione}

\subsection{L'iterazione zero}

All'iterazione $j=1$ il problema~\eqref{eq:lmpc} usa $\CS^0$ e $P^0$, costruiti
con la sola iterazione $0$. Serve quindi che esista almeno una prima
traiettoria buona.

\begin{ipotesi}[Assumption~1 del paper]\label{ip:iterazione-zero}
È disponibile una traiettoria $\vett{x}^0$ dell'iterazione $0$, con ingressi
$\vett{u}^0$, che parte da $x_S$, rispetta i vincoli e converge a $x_F$. Di
conseguenza $0\in\Mc^0$ e $\CS^0$ non è vuoto.
\end{ipotesi}

La traiettoria iniziale non deve essere buona, deve solo essere
\emph{ammissibile}: può venire da un controllore semplice e conservativo, da
una manovra lenta, perfino da una dimostrazione manuale. Da lì in poi è il
LMPC a migliorare.

\subsection{Lo schema complessivo}

\begin{figure}[h]
\centering
\begin{tikzpicture}[>=Stealth, font=\small, node distance=1.2cm,
   blocco/.style={draw, rounded corners, align=center, minimum height=1.1cm, fill=gray!6}]
  \node[blocco, minimum width=3.4cm] (lmpc) at (0,0) {LMPC\\ risolve il QP~\eqref{eq:lmpc-unico}};
  \node[blocco, minimum width=3.4cm] (sys) at (6,0) {sistema\\ $x_{t+1}=Ax_t+Bu_t$};
  \node[blocco, minimum width=4.6cm] (mem) at (0,2.6)
     {dati delle iterazioni $0,\dots,j-1$\\ $\vett{D}^{j-1}$ (stati), $\vett{c}^{j-1}$ (costi)};
  \draw[->, thick] (lmpc) -- node[above]{$u^j_t$} (sys);
  \draw[->, thick] (sys.east) -- ++(0.7,0) |- ($(lmpc.south)+(0,-0.7)$)
      node[pos=0.75, below]{$x^j_t$} -- (lmpc.south);
  \draw[->, thick] (mem) -- node[right]{$\CS^{j-1}$, $P^{j-1}$} (lmpc);
  \draw[->, dashed] (sys.north) |- (mem.east);
  \node[right, align=left, font=\footnotesize] at ($(sys.north)+(0.1,1.2)$)
      {a fine iterazione:\\ si aggiunge la traiettoria $j$\\ e i suoi costi};
\end{tikzpicture}
\caption{Schema del LMPC. Dentro un'iterazione è un MPC con insieme e costo
terminali fissi. Tra un'iterazione e l'altra i dati vengono aggiornati.}
\label{fig:schema}
\end{figure}

\medskip
\noindent\fbox{\begin{minipage}{0.96\textwidth}
\smallskip
\textbf{Algoritmo LMPC (sistemi lineari).}
\begin{enumerate}
  \item \emph{Inizializzazione.} Si esegue l'iterazione $0$ con un controllore
        qualsiasi che porti da $x_S$ a $x_F$ nel rispetto dei vincoli. Si
        memorizzano stati e ingressi e si calcolano i costi realizzati con
        la~\eqref{eq:ricorsione-costi}.
  \item \emph{Iterazione $j=1,2,\dots$}: si pone $t=0$, $x^j_0=x_S$; finché
        l'iterazione non è conclusa:
  \begin{enumerate}
    \item si misura $x^j_t$;
    \item si risolve il QP~\eqref{eq:lmpc-unico} con i dati delle iterazioni
          $0,\dots,j-1$;
    \item si applica $u^j_t=u^{*,j}_{t|t}$ e si pone $t\leftarrow t+1$.
  \end{enumerate}
  \item \emph{Aggiornamento.} Si aggiungono a $\vett{D}$ gli stati
        dell'iterazione $j$ e a $\vett{c}$ i loro costi realizzati, calcolati
        a ritroso. Si passa all'iterazione $j+1$.
\end{enumerate}
\smallskip
\end{minipage}}
\medskip

% ===========================================================================
\section{Fattibilità ricorsiva e stabilità}\label{sec:stabilita}
% ===========================================================================

\begin{teorema}[Teorema~1 del paper]\label{teo:1}
Si consideri il sistema~\eqref{eq:sistema} controllato dal LMPC~\eqref{eq:lmpc},
\eqref{eq:legge}, con $\CS^j$ definito da~\eqref{eq:CS}, e valga
l'ipotesi~\ref{ip:iterazione-zero}. Allora, per ogni iterazione $j\ge 1$:
\begin{enumerate}
  \item il problema~\eqref{eq:lmpc} ha soluzione per ogni $t\in\Zn$
        (e quindi i vincoli non sono mai violati);
  \item l'equilibrio $x_F$ è asintoticamente stabile per l'anello
        chiuso~\eqref{eq:sistema}, \eqref{eq:legge}.
\end{enumerate}
\end{teorema}

Il paper scrive che ``la dimostrazione segue da argomenti standard di MPC''.
Alla luce della sezione~\ref{sec:richiami} il motivo è chiaro: è il
teorema~\ref{teo:mpc} con $(\Xf,V_f)=(\CS^{j-1},P^{j-1})$, e il
lemma~\ref{lem:decrescita} garantisce che questa coppia soddisfa
l'ipotesi~\ref{ip:terminali}. Conviene comunque seguire i passi del paper, che
rendono esplicito il ruolo dei dati.

\paragraph{Passo 1: fattibilità all'istante iniziale (\eqp{21}--\eqp{22}).}
Per la~\eqref{eq:CS-monotono}, $\CS^0\subseteq\CS^{j-1}$ per ogni $j\ge 1$, e
$\CS^0$ non è vuoto per l'ipotesi~\ref{ip:iterazione-zero}. All'istante $t=0$
dell'iterazione $j$ lo stato è $x^j_0=x_S=x^0_0$. I primi $N$ passi
dell'iterazione zero,
\[
  [x^0_0,\,x^0_1,\,\dots,\,x^0_N],\qquad [u^0_0,\,u^0_1,\,\dots,\,u^0_{N-1}],
\]
sono una soluzione ammissibile di~\eqref{eq:lmpc}: partono dallo stato giusto,
rispettano dinamica e vincoli (sono stati realizzati) e terminano in
$x^0_N\in\SSet^0\subseteq\CS^0\subseteq\CS^{j-1}$.

\paragraph{Passo 2: dalla fattibilità in $t$ a quella in $t+1$ (\eqp{23}--\eqp{28}).}
Sia il problema risolto all'istante $t$, con soluzione ottima
$\vett{u}^{*,j}_{t:t+N|t}$, $\vett{x}^{*,j}_{t:t+N|t}$. Il vincolo terminale dà
$x^{*,j}_{t+N|t}\in\CS^{j-1}$; siano $\lambda^{*,i}_s$ i pesi che realizzano
$P^{j-1}(x^{*,j}_{t+N|t})$ (\eqp{24}). Con questi pesi si costruiscono
(\eqp{25}--\eqp{26})
\[
  \bar u = \sum_{i,s}\lambda^{*,i}_s u^i_s\in\Uc,\qquad
  \bar x = Ax^{*,j}_{t+N|t}+B\bar u = \sum_{i,s}\lambda^{*,i}_s x^i_{s+1}\in\CS^{j-1}.
\]
Poiché il modello di predizione coincide con il sistema, lo stato realizzato
all'istante successivo è quello predetto (\eqp{27}):
\begin{equation}\label{eq:modello-esatto}
  x^j_{t+1} = x^{*,j}_{t+1|t}.
\end{equation}
All'istante $t+1$ la sequenza traslata (\eqp{28})
\[
  [u^{*,j}_{t+1|t},\,\dots,\,u^{*,j}_{t+N-1|t},\ \bar u],\qquad
  [x^{*,j}_{t+1|t},\,\dots,\,x^{*,j}_{t+N|t},\ \bar x]
\]
parte dallo stato corretto, rispetta dinamica e vincoli e termina in
$\CS^{j-1}$: è ammissibile (figura~\ref{fig:traslazione}). Per induzione su
$t$, il problema ha soluzione a ogni istante.

\begin{figure}[h]
\centering
\begin{tikzpicture}[>=Stealth, font=\small, scale=1.1]
  \fill[blue!10] (4,3) -- (5.5,3.6) -- (5.2,2.2) -- cycle;
  \fill[teal!12] (5,2.4) -- (6.4,2.9) -- (6,1.5) -- cycle;
  \foreach \a/\b in {{(4,3)}/{(5,2.4)}, {(5.5,3.6)}/{(6.4,2.9)}, {(5.2,2.2)}/{(6,1.5)}}
     \draw[->, gray] \a -- \b;
  \foreach \p in {(4,3),(5.5,3.6),(5.2,2.2)} \fill[blue!60!black] \p circle (1.8pt);
  \foreach \p in {(5,2.4),(6.4,2.9),(6,1.5)} \fill[teal!70!black] \p circle (1.8pt);
  \node[blue!60!black, above left] at (4,3) {$x^{i}_{s}$};
  \node[teal!70!black, below right] at (6,1.5) {$x^{i}_{s+1}$};
  \draw[thick, blue] (0,1) -- (1.4,2.0) -- (2.8,2.6) -- (4.93,2.86);
  \foreach \p in {(1.4,2.0),(2.8,2.6)} \fill[blue] \p circle (1.8pt);
  \fill (0,1) circle (2pt) node[below]{$x^j_t$};
  \fill[blue] (4.93,2.86) circle (2.2pt);
  \draw[blue, thin] (4.93,2.86) -- (4.5,3.75) node[above]{$x^{*,j}_{t+N|t}$};
  \draw[thick, dashed, red!80!black] (1.4,2.0) -- (2.8,2.6) -- (4.93,2.86) -- (5.82,2.19);
  \fill[red!80!black] (5.82,2.19) circle (2.2pt) node[below left]{$\bar x$};
  \node[red!80!black, font=\footnotesize] at (5.55,2.55) {$\bar u$};
  \node[below right] at (1.4,2.0) {$x^j_{t+1}$};
\end{tikzpicture}
\caption{Argomento della traslazione. Lo stato finale predetto all'istante $t$
(blu) è combinazione convessa di alcuni stati memorizzati (triangolo blu).
Applicando la stessa combinazione dei loro ingressi si arriva in $\bar x$, la
stessa combinazione dei loro successori (triangolo verde), che è ancora in
$\CS^{j-1}$. La traiettoria rossa tratteggiata è la candidata ammissibile
all'istante $t+1$.}
\label{fig:traslazione}
\end{figure}

\paragraph{Passo 3: il costo ottimo come funzione di Lyapunov (\eqp{29}--\eqp{32}).}
All'interno dell'iterazione $j$ il problema~\eqref{eq:lmpc} non dipende da
$t$ (dinamica, vincoli, $\CS^{j-1}$ e $P^{j-1}$ sono fissi): il suo valore ottimo
è una funzione del solo stato, che si indica con $\JL(\cdot)$. Si vuole
mostrare che è una funzione di Lyapunov per $x_F$.

\begin{itemize}
  \item \emph{Positività.} Per $x\ne x_F$ il costo ottimo contiene almeno il
        termine $h(x,u_{t|t})>0$; in $x_F$ la soluzione $u=0$ mantiene lo
        stato in $x_F\in\CS^{j-1}$ con costo nullo, quindi $\JL(x_F)=0$.
  \item \emph{Continuità.} Si dimostra come in Mayne et al.~(2000); il paper
        non entra nei dettagli.
  \item \emph{Decrescita.} Si scrive il costo ottimo in $t$ separando il primo
        termine e sviluppando $P^{j-1}$ con la~\eqref{eq:P-esplicita}
        (\eqp{30}); con la disuguaglianza di Jensen e la definizione di
        $P^{j-1}$ come minimo, esattamente come nel lemma~\ref{lem:decrescita},
        si ottiene (\eqp{31})
        \[
        \begin{aligned}
          \JL(x^j_t) &\ge h(x^j_t,u^j_t)
          + \underbrace{\sum_{k=t+1}^{t+N-1}h(x^{*,j}_{k|t},u^{*,j}_{k|t})
          + h(x^{*,j}_{t+N|t},\bar u) + P^{j-1}(\bar x)}_{\text{costo della candidata in } t+1}\\
          &\ge h(x^j_t,u^j_t) + \JL(x^j_{t+1}).
        \end{aligned}
        \]
        L'ultima disuguaglianza vale perché il costo ottimo in $t+1$ non supera
        quello della candidata ammissibile del passo 2.
\end{itemize}
Si ottiene così (\eqp{32})
\begin{equation}\label{eq:lyapunov-lmpc}
  \JL(x^j_{t+1}) - \JL(x^j_t) \le -h(x^j_t,u^j_t) < 0
  \qquad \forall x^j_t\ne x_F,
\end{equation}
e il teorema di Lyapunov dà la stabilità asintotica di $x_F$.

\begin{osservazione}[Regione di attrazione]
Il LMPC all'iterazione $j$ funziona da ogni stato da cui sia possibile entrare
in $\CS^{j-1}$ in $N$ passi rispettando i vincoli. Nel linguaggio del libro è
l'\emph{insieme controllabile in $N$ passi} $K_N(\CS^{j-1})$
(sezione~\ref{sec:libro-insiemi}). Poiché $\CS^{j-1}$ cresce con
$j$, cresce anche questa regione: le iterazioni non migliorano solo il costo,
ma anche l'insieme degli stati da cui il controllore è in grado di operare.
\end{osservazione}

% ===========================================================================
\section{Miglioramento e ottimalità}\label{sec:convergenza}
% ===========================================================================

\subsection{Il costo non peggiora da un'iterazione all'altra}

\begin{teorema}[Teorema~2 del paper]\label{teo:2}
Nelle ipotesi del teorema~\ref{teo:1}, il costo di ogni iterazione non
supera quello della precedente:
\[
  J^{j}_{0\to\infty}(x_S)\le J^{j-1}_{0\to\infty}(x_S)\qquad\forall j\ge 1.
\]
\end{teorema}

Il paper rimanda la dimostrazione al lavoro del 2016. Con gli strumenti di
questo documento la si ricostruisce in due passaggi.

\paragraph{Prima disuguaglianza: il costo predetto in $x_S$ non supera il costo dell'iterazione precedente.}
L'iterazione $j-1$ è riuscita (teorema~\ref{teo:1} o ipotesi~\ref{ip:iterazione-zero}
se $j=1$), quindi i suoi dati sono in $\CS^{j-1}$. I suoi primi $N$ passi sono
una soluzione ammissibile di~\eqref{eq:lmpc} in $x_S$ (come nel passo 1 del
teorema~\ref{teo:1}). Il loro costo, usando la proposizione~\ref{prop:P}(c), è
\[
  \sum_{k=0}^{N-1}h(x^{j-1}_k,u^{j-1}_k) + P^{j-1}(x^{j-1}_N)
  \le \sum_{k=0}^{N-1}h(x^{j-1}_k,u^{j-1}_k) + J^{j-1}_{N\to\infty}(x^{j-1}_N)
  = J^{j-1}_{0\to\infty}(x_S).
\]
Il costo ottimo non supera quello di una soluzione ammissibile:
$\JL(x_S)\le J^{j-1}_{0\to\infty}(x_S)$.

\paragraph{Seconda disuguaglianza: il costo realizzato non supera il costo predetto.}
Sommando la~\eqref{eq:lyapunov-lmpc} da $t=0$ a $t=T$ i termini si
cancellano a due a due (somma telescopica):
\[
  \sum_{t=0}^{T}h(x^j_t,u^j_t) \le \JL(x^j_0) - \JL(x^j_{T+1}) \le \JL(x_S),
\]
perché $\JL\ge 0$. Per $T\to\infty$: $J^j_{0\to\infty}(x_S)\le\JL(x_S)$.

\medskip
\noindent
Mettendo insieme:
\begin{equation}\label{eq:catena}
  \Jstar(x_S)\ \le\ J^j_{0\to\infty}(x_S)\ \le\ \JL(x_S)\ \le\ J^{j-1}_{0\to\infty}(x_S).
\end{equation}
La prima disuguaglianza vale perché la traiettoria dell'iterazione $j$ è
ammissibile per il problema a orizzonte infinito. La successione dei costi di
iterazione è quindi \emph{non crescente e limitata inferiormente}: converge
sempre a un limite. Il teorema successivo dice qual è.

\subsection{La traiettoria limite è ottima}

\begin{teorema}[Teorema~3 del paper]\label{teo:3}
Nelle ipotesi del teorema~\ref{teo:1}, con $N>1$, si supponga che le
traiettorie in anello chiuso convergano, per $j\to\infty$, a una traiettoria
stazionaria $\vett{x}^\infty=\lim_j\vett{x}^j$ con ingressi
$\vett{u}^\infty=\lim_j\vett{u}^j$. Allora $(\vett{x}^\infty,\vett{u}^\infty)$
è una soluzione ottima globale del problema a orizzonte
infinito~\eqref{eq:ottimo-infinito}: $\vett{x}^\infty=\vett{x}^*$,
$\vett{u}^\infty=\vett{u}^*$.
\end{teorema}

Anche questa dimostrazione è rimandata al lavoro del 2016; qui se ne dà
l'idea, senza pretesa di completezza.
\begin{itemize}
  \item Quando le iterazioni non cambiano più, la traiettoria stazionaria fa
        parte dei dati, e lungo di essa il costo terminale coincide con il
        costo realmente pagato. A ogni istante il LMPC potrebbe scegliere di
        deviarne per $N$ passi e poi rientrare in $\CS$; se non lo fa, è perché
        \emph{nessuna modifica di $N$ ingressi consecutivi} migliora il costo.
        La traiettoria stazionaria è quindi un punto in cui non esistono
        miglioramenti ``locali'' (nel senso di modifiche su finestre di $N$
        passi).
  \item Il problema a orizzonte infinito~\eqref{eq:ottimo-infinito}, per un
        sistema lineare con vincoli convessi e costo convesso, è un problema
        \emph{convesso}. In un problema convesso un punto che non si può
        migliorare localmente è un ottimo globale
        (sezione~\ref{sec:ottimizzazione}).
\end{itemize}
Per sistemi non lineari (il lavoro del 2016) lo stesso ragionamento dà solo
un ottimo locale; è la convessità, cioè ancora una volta la linearità del
sistema, a renderlo globale. (L'introduzione della sezione 3.3 del paper parla
di ottimo ``locale'', ma l'enunciato del teorema~3 è di ottimalità globale,
coerentemente con questa osservazione.)

Due cautele nella lettura del teorema~\ref{teo:3}:
\begin{itemize}
  \item la convergenza delle traiettorie è un'\emph{ipotesi}, non una tesi: il
        teorema~\ref{teo:2} garantisce che i \emph{costi} convergono, ma non che
        le traiettorie convergano;
  \item l'ipotesi $N>1$ è richiesta dalla dimostrazione del 2016; il paper non
        commenta che cosa possa succedere con $N=1$.
\end{itemize}

\begin{figure}[h]
\centering
\begin{tikzpicture}[>=Stealth, font=\small, x=0.9cm, y=1.3cm]
  \draw[->] (-0.2,0) -- (8,0) node[right]{iterazione $j$};
  \draw[->] (0,-0.1) -- (0,3.4) node[left]{costo di iterazione};
  \draw[dashed, red!70!black] (0,0.8) -- (7.6,0.8) node[right]{$\Jstar(x_S)$};
  \foreach \j in {0,...,7}{
     \fill[blue!70!black] ({\j},{0.8+2.2*exp(-1.6*\j)}) circle (2.2pt);
     \draw (\j,0.05) -- (\j,-0.05) node[below]{$\j$};
  }
  \draw[blue!70!black, thin] plot[domain=0:7, samples=40] (\x,{0.8+2.2*exp(-1.6*\x)});
  \node[blue!70!black, right] at (0.2,3.0) {$J^j_{0\to\infty}(x_S)$};
\end{tikzpicture}
\caption{Andamento qualitativo del costo di iterazione previsto dai
teoremi~\ref{teo:2} e~\ref{teo:3}: non cresce mai e, se le traiettorie
convergono, tende al costo ottimo a orizzonte infinito. Nell'esempio del paper
quasi tutto il miglioramento avviene alla prima iterazione.}
\label{fig:costo-iterazioni}
\end{figure}

% ===========================================================================
\section{L'esempio del paper: LQR vincolato}\label{sec:esempio}
% ===========================================================================

\paragraph{Il problema di prova.}
Il paper verifica il LMPC su un regolatore lineare quadratico vincolato a
orizzonte infinito (\eqp{33}):
\[
  \min_{u_0,u_1,\dots}\sum_{k=0}^{\infty}\big(\|x_k\|_2^2+\|u_k\|_2^2\big),
  \qquad
  x_{k+1}=\begin{bmatrix}1&1\\0&1\end{bmatrix}x_k+\begin{bmatrix}0\\1\end{bmatrix}u_k,
\]
con $x_0=x_S$, limiti $-4\le x_k\le 4$ su ciascuna componente e $-1\le u_k\le 1$.
Il sistema è un \emph{doppio integratore} discreto: la prima componente dello
stato è una posizione, la seconda una velocità, l'ingresso un'accelerazione.
L'equilibrio è $x_F=0$, il costo è quadratico e definito positivo, i vincoli
sono scatole: tutte le ipotesi del paper sono soddisfatte. La scelta è
didattica per un motivo preciso: per questo problema la soluzione ottima
esatta a orizzonte infinito si può calcolare (con l'algoritmo di Borrelli,
2003), e quindi si può verificare se il LMPC la trova davvero.

\paragraph{Implementazione.}
Il LMPC è scritto come QP (sezione~\ref{sec:qp}) con YALMIP e risolto con
\texttt{quadprog}. Per avere iterazioni di durata finita, il paper considera
conclusa l'iterazione $j$ al primo istante in cui il costo ottimo predetto
scende sotto una soglia piccolissima (\eqp{34}):
\[
  t_j = \min\big\{t\in\Zn : \JL(x^j_t)\le\varepsilon\big\},
\]
cioè quando lo stato è, a tutti gli effetti, arrivato all'equilibrio. Le
somme della funzione baricentrica sono troncate a $t_j$ come
in~\eqp{17}.

\paragraph{Risultati (lettura qualitativa).}
\begin{itemize}
  \item \emph{Miglioramento monotono} (Table~1 del paper). Il costo di
        iterazione non cresce mai, come previsto dal teorema~\ref{teo:2}. Il
        salto più grande avviene alla prima iterazione (la traiettoria
        iniziale è volutamente non ottima); le iterazioni successive limano
        differenze sempre più piccole, fino a stabilizzarsi.
  \item \emph{Convergenza all'ottimo vero} (equazioni \eqp{36}--\eqp{39}). La
        traiettoria stazionaria del LMPC coincide con quella ottima del CLQR
        calcolata in modo esatto, a meno di errori dell'ordine della
        precisione numerica, sia come stati sia come costo. È la verifica
        sperimentale del teorema~\ref{teo:3}.
  \item \emph{Ruolo dell'orizzonte} (Table~2). Il risultato vale per diversi
        stati iniziali e diversi orizzonti, anche cortissimi ($N=2$): un
        orizzonte corto, accompagnato da un costo terminale imparato, basta a
        raggiungere l'ottimo a orizzonte \emph{infinito}. Quello che cambia è
        la velocità: con orizzonte più lungo il controllore ha più libertà di
        esplorare lo spazio di stato in ogni iterazione, e converge in molte
        meno iterazioni (decine con $N=2$, una decina con $N=4$).
  \item \emph{Costo computazionale}. Sullo stesso problema, il LMPC convesso
        converge in decine di secondi, quello del 2016 basato su MIP in ore.
        È la conferma del contributo principale del paper.
\end{itemize}

% ===========================================================================
\section{Collegamenti con il libro di Borrelli, Bemporad e Morari}\label{sec:libro}
% ===========================================================================

Il paper occupa sei pagine e per quasi tutti i risultati di base rimanda ad
``argomenti standard di MPC'' (Mayne et al., 2000) e a Borrelli (2003). Quegli
argomenti sono esposti in modo sistematico nel libro \emph{Predictive Control
for Linear and Hybrid Systems} di F.~Borrelli, A.~Bemporad e M.~Morari
(2017), indicato qui come [BBM]. Leggere il paper con il libro accanto mostra
una cosa importante: \textbf{il LMPC non introduce una nuova teoria della
stabilità}. Riusa il teorema di stabilità dell'MPC del libro, e sostituisce
gli ingredienti terminali, che nel libro sono progettati con il regolatore
LQR e con operazioni sui poliedri, con ingredienti ricavati dai dati. Anche lo
strumento con cui li ricava, l'interpolazione baricentrica, è nel libro, nel
capitolo~13, usato però per uno scopo diverso.

La tabella~\ref{tab:mappa} riassume i collegamenti; le sottosezioni li
discutono uno per uno.

\begin{table}[h]
\centering
\caption{Dove si trovano nel libro i concetti usati dal paper.}\label{tab:mappa}
\small
\begin{tabular}{p{4.1cm}p{3.6cm}p{6.6cm}}
\toprule
Nel paper & Nel libro & Che cosa dice il libro \\
\midrule
Problema a orizzonte infinito \eqp{5}, CLQR \eqp{33}
  & Sez.~7.4, 11.3.4--11.3.5
  & equazione di Bellman, soluzione esatta del CLQR affine a tratti
    (Teorema~11.4) e algoritmo per calcolarla, usato dal paper come
    riferimento \\
Safe set, invarianza
  & Def.~10.9--10.13, Teoremi~11.2 e~12.1
  & insiemi controllabili e invarianti, insiemi ammissibili dell'MPC,
    fattibilità ricorsiva \\
Inviluppo convesso \eqp{10}
  & Sez.~4.2, 4.4.2
  & rappresentazioni V e H dei politopi, costo del cambio di
    rappresentazione \\
Teorema~1 (stabilità)
  & Teorema~12.2, Remark~12.2--12.3
  & ipotesi (A0)--(A3) su insieme e costo terminali, funzione di Lyapunov
    di controllo \\
Funzione baricentrica \eqp{14}--\eqp{15}
  & Cap.~13 (Def.~13.1, Lemmi~13.1--13.2, Teorema~13.2)
  & interpolazione baricentrica, proprietà dell'inviluppo convesso,
    sovrastima del costo ottimo \\
Convessità di $P^j$, Remark~4
  & Teorema~6.5, Corollario~5.1, Teorema~6.7
  & la funzione valore di un programma multiparametrico convesso è
    convessa; per un mp-LP è anche continua e affine a tratti \\
LMPC non lineare del 2016 (MIP)
  & Sez.~2.4, 16.4, 17.4
  & vincoli logici come variabili binarie, MIQP, branch and bound \\
Esempio: doppio integratore
  & Esempi~11.1, 11.2, 12.1
  & lo stesso sistema, usato per insiemi ammissibili, CLQR esplicito e
    perdita di fattibilità del receding horizon \\
Limiti: disturbi, incertezze
  & Sez.~10.3--10.4, Cap.~15
  & insiemi invarianti robusti, MPC robusto \\
\bottomrule
\end{tabular}
\end{table}

\paragraph{Corrispondenza di notazione.}
Il libro chiama $q(x,u)$ il costo di stadio (qui $h$), $p(x)$ il costo
terminale (qui $V_f$, e nel LMPC $P^{j-1}$) e $\Xf$ l'insieme terminale (nel
LMPC $\CS^{j-1}$). Il libro usa la lettera $P$ per la \emph{matrice} del costo
terminale quadratico $x^\top Px$: non va confusa con la \emph{funzione} $P^j$
del paper. Infine il libro pone l'equilibrio nell'origine, mentre il paper lo
chiama $x_F$; poiché $x_F = Ax_F$, il cambio di variabile $x\mapsto x-x_F$
trasforma il sistema del paper in uno con equilibrio nell'origine, con la
stessa dinamica, e i risultati si trasferiscono senza modifiche.

\subsection{Il Teorema~1 del paper è il Teorema~12.2 del libro}\label{sec:libro-teorema}

Il Teorema~12.2 del libro è il risultato di riferimento sulla fattibilità e
stabilità dell'MPC: il libro stesso, nella sezione~12.8, lo indica come ``il
risultato principale'' su fattibilità e stabilità e ne fa risalire l'idea a
Keerthi e Gilbert. Le sue ipotesi, nella notazione di questo documento, e il modo in
cui il LMPC le soddisfa sono:
\begin{itemize}
  \item \textbf{(A0)} costo di stadio e costo terminale continui e definiti
        positivi. Per $h$ è l'ipotesi~\eqp{6}. Per $P^{j-1}$: è continua su
        $\CS^{j-1}$ (sezione~\ref{sec:libro-mp}), si annulla in $x_F$ e per
        $x\ne x_F$ vale $P^{j-1}(x)\ge\Jstar(x)>0$
        (proposizione~\ref{prop:P}(b)).
  \item \textbf{(A1)} $\Xc$, $\Uc$, $\Xf$ chiusi e contenenti l'equilibrio
        nel loro \emph{interno}: su questo punto si veda l'osservazione
        qui sotto.
  \item \textbf{(A2)} $\Xf$ controllo-invariante e contenuto in $\Xc$: è la
        proposizione~\ref{prop:CS-invariante}.
  \item \textbf{(A3)}
        $\displaystyle\min_{v\in\Uc,\,Ax+Bv\in\Xf}\big(-p(x)+q(x,v)+p(Ax+Bv)\big)\le 0$
        per ogni $x\in\Xf$: è il lemma~\ref{lem:decrescita}, con
        $v=\bar u$.
\end{itemize}
La dimostrazione del libro (equazioni~(12.19)--(12.20) di [BBM]) è la stessa
del paper (\eqp{28}--\eqp{32}): sequenza traslata più un ingresso finale
fornito da (A3), confronto con l'ottimo, decrescita del costo ottimo. L'unica
parte che il paper deve scrivere esplicitamente è la verifica di (A3), che nel
libro è garantita in fase di progetto (con il costo LQR, equazioni~(12.24)--(12.26)
di [BBM], vale con il segno di uguaglianza) e che nel paper discende dalla
disuguaglianza di Jensen applicata ai dati. In questo senso il contributo
teorico del paper sta tutto nella sezione~\ref{sec:costo-terminale}: mostrare
che $(\CS^{j-1},P^{j-1})$ soddisfa (A2) e (A3).

Il libro contiene anche, nella sezione~7.4.2 (Teorema~7.1), la versione più
semplice dello stesso ragionamento, con vincolo terminale $x_N=0$: la sequenza
ottima traslata e completata con ingresso nullo è ammissibile al passo
successivo, e il costo ottimo decresce del costo di stadio. È utile
leggerla prima del Teorema~12.2, perché isola l'idea della traslazione senza
gli ingredienti terminali.

\begin{osservazione}[L'ipotesi (A1) e l'equilibrio sul bordo di $\CS^{j-1}$]
Nel libro (A1) serve a dimostrare la continuità del costo ottimo
nell'equilibrio, tramite la disuguaglianza $J^*_0(x)\le p(x)$ su $\Xf$
(equazione~(12.23) di [BBM]). Nel LMPC l'equilibrio $x_F$ è l'ultimo punto di
tutte le traiettorie memorizzate e di solito sta sul \emph{bordo}
dell'inviluppo convesso, spesso su un vertice (si veda la
figura~\ref{fig:safe-set}), non al suo interno. La disuguaglianza
$\JL(x)\le P^{j-1}(x)$ su $\CS^{j-1}$ si ottiene comunque con lo stesso
argomento del libro (iterando il lemma~\ref{lem:decrescita}), ma la stabilità
va intesa rispetto all'insieme degli stati ammissibili, come nel teorema di
Lyapunov del libro (Teorema~7.2, enunciato su un insieme che contiene
l'origine, non necessariamente nel suo interno). Il paper non entra in questo
dettaglio e rimanda a Mayne et al.\ per la continuità.
\end{osservazione}

\begin{osservazione}[Il problema che il LMPC risolve, nelle parole del libro]
Nel Remark~12.2 il libro osserva che la procedura del Teorema~12.2 è in
generale \emph{conservativa}: introdurre un insieme terminale artificiale
riduce la regione di attrazione e può peggiorare le prestazioni. Con la
scelta standard ($\Xf$ = massimo insieme invariante del sistema in anello
chiuso con il guadagno LQR, costo terminale $x^\top P_\infty x$) l'insieme
terminale è un intorno dell'equilibrio in cui i vincoli non sono attivi, di
solito piccolo. Il LMPC risponde proprio a questa osservazione: l'insieme
terminale non è un intorno piccolo costruito attorno all'equilibrio, ma
l'insieme di tutti gli stati da cui si sa già, per esperienza, come arrivare
all'equilibrio; e cresce con le iterazioni.
\end{osservazione}

\subsection{Insiemi controllabili, invarianti e regione di attrazione}\label{sec:libro-insiemi}

Il capitolo~10 del libro dà il linguaggio preciso per parlare della regione in
cui il LMPC funziona. Le definizioni che servono sono:
\begin{itemize}
  \item l'\emph{insieme predecessore} \libro{Def.~10.2}
        $\mathrm{Pre}(S)=\{x: \exists u\in\Uc,\ Ax+Bu\in S\}$, gli stati che
        possono essere portati in $S$ in un passo;
  \item l'\emph{insieme controllabile in $N$ passi} \libro{Def.~10.5},
        $K_N(S)$, definito ricorsivamente da $K_0(S)=S$ e
        $K_i(S)=\mathrm{Pre}(K_{i-1}(S))\cap\Xc$: gli stati che possono essere
        portati in $S$ in $N$ passi rispettando i vincoli;
  \item il \emph{massimo insieme controllo-invariante} $\mathcal{C}_\infty$
        \libro{Def.~10.10}, il più grande insieme in cui si può tenere lo
        stato per sempre rispettando i vincoli;
  \item per il CLQR, l'insieme $\Xc_\infty$ degli stati da cui il problema
        a orizzonte infinito ha costo finito \libro{eq.~(11.44)}, che il
        libro identifica con il massimo insieme stabilizzabile verso
        l'equilibrio. È il dominio di $\Jstar$.
\end{itemize}
Il Teorema~11.2 del libro dice che, se l'insieme terminale $\Xf$ è
controllo-invariante, l'insieme degli stati da cui il problema MPC ha
soluzione è $K_N(\Xf)$, che è a sua volta controllo-invariante e contenuto in
$\mathcal{C}_\infty$. Il Teorema~12.1 ne ricava la fattibilità ricorsiva,
\emph{qualunque sia la funzione di costo}. Applicati al LMPC:
\begin{itemize}
  \item all'iterazione $j$ il LMPC ha soluzione esattamente in
        $K_N(\CS^{j-1})$, che è controllo-invariante;
  \item la fattibilità ricorsiva dipende solo da $\CS^{j-1}$, non da
        $P^{j-1}$: il costo terminale serve per la stabilità e per il
        miglioramento, non per la fattibilità;
  \item poiché ogni punto di $\CS^{j-1}$ ha costo $\Jstar$ finito
        (proposizione~\ref{prop:P}(b)), vale la catena di inclusioni
        \begin{equation}\label{eq:inclusioni}
          \{x_F\}\subseteq\CS^0\subseteq\CS^{j-1}\subseteq K_N(\CS^{j-1})
          \subseteq\Xc_\infty\subseteq\mathcal{C}_\infty\subseteq\Xc,
        \end{equation}
        illustrata nella figura~\ref{fig:inclusioni}.
\end{itemize}

\begin{figure}[h]
\centering
\begin{tikzpicture}[font=\small, scale=1.05]
  \draw[gray] (0,0) rectangle (9,5);
  \node[gray, below left] at (9,5) {$\Xc$};
  \draw[fill=gray!10] (4.5,2.5) ellipse (4.2 and 2.2);
  \node at (4.5,4.4) {$\mathcal{C}_\infty$};
  \draw[fill=orange!12] (4.8,2.4) ellipse (3.4 and 1.75);
  \node at (4.8,3.85) {$\Xc_\infty$};
  \draw[fill=green!12] (5.2,2.3) ellipse (2.5 and 1.3);
  \node at (5.0,3.25) {$K_N(\CS^{j-1})$};
  \fill[blue!18] (4.2,2.2) -- (5.0,2.9) -- (6.6,2.6) -- (7.0,1.9) -- (5.6,1.5) -- cycle;
  \draw[blue!60] (4.2,2.2) -- (5.0,2.9) -- (6.6,2.6) -- (7.0,1.9) -- (5.6,1.5) -- cycle;
  \fill[blue!35] (4.2,2.2) -- (5.0,2.5) -- (5.9,2.0) -- (5.6,1.5) -- cycle;
  \node[font=\footnotesize] at (5.15,2.05) {$\CS^0$};
  \node[font=\footnotesize] at (6.35,2.25) {$\CS^{j-1}$};
  \fill (4.2,2.2) circle (1.6pt) node[left]{$x_S$};
  \fill (5.6,1.5) circle (1.6pt) node[below]{$x_F$};
\end{tikzpicture}
\caption{Gli insiemi della catena~\eqref{eq:inclusioni} (schematico). Il
safe set convesso cresce con le iterazioni dentro $\Xc_\infty$, il dominio
della funzione valore ottima $\Jstar$; la regione in cui il LMPC ha soluzione è l'insieme degli stati
da cui lo si raggiunge in $N$ passi.}
\label{fig:inclusioni}
\end{figure}

Il confronto con il modo in cui il libro calcola questi insiemi è istruttivo.
Gli algoritmi~10.2 e~10.3 di [BBM] calcolano $\mathcal{C}_\infty$ e gli insiemi
stabilizzabili iterando l'operatore $\mathrm{Pre}$, cioè proiezioni di poliedri
in rappresentazione H; in generale non è garantito che terminino in un numero
finito di passi, e il costo cresce rapidamente con la dimensione dello stato.
Il LMPC ottiene invece, \emph{senza alcuna operazione sugli insiemi}, un
sottoinsieme controllo-invariante di $\Xc_\infty$ in rappresentazione V: i
dati stessi ne sono i generatori. È un'approssimazione interna, che dipende
da quali regioni dello spazio di stato le iterazioni hanno visitato, ma è
disponibile gratuitamente e si allarga da sola.

\subsection{La funzione baricentrica e il capitolo~13}\label{sec:libro-baricentrica}

Il paper attribuisce la funzione baricentrica a Jones e Morari (2010). Il
capitolo~13 del libro, scritto da C.~N.~Jones, espone la stessa teoria nel suo
contesto originale: l'\emph{MPC approssimato}. Lì il problema è un altro: si ha
già un MPC con buone proprietà, ma risolverlo online, o memorizzarne la
soluzione esplicita, costa troppo. Si calcola allora la soluzione ottima
$z^*(v)$ in alcuni punti campione $v\in V$ dello spazio di stato e la si
interpola nei punti intermedi.

\paragraph{Gli strumenti del capitolo 13.}
\begin{itemize}
  \item \libro{Def.~13.1}: un insieme di funzioni $w_v(x)$ è
        \emph{baricentrico} su $\Conv(V)$ se $w_v(x)\ge 0$ (positività),
        $\sum_v w_v(x)=1$ (partizione dell'unità) e $\sum_v v\,w_v(x)=x$
        (precisione lineare). Sono esattamente i vincoli sui pesi $\lambda$
        nel problema~\eqref{eq:LP} (\eqp{15b}--\eqp{15c}).
  \item \libro{eq.~(13.18)}: la soluzione interpolata è
        $\tilde z(x)=\sum_v w_v(x)\,z^*(v)$.
  \item \libro{Lemma~13.1} (\emph{proprietà dell'inviluppo convesso}): la
        soluzione interpolata è ammissibile in tutto $\Conv(V)$. Il motivo è
        che i vincoli sono convessi congiuntamente in soluzione e stato: è lo
        stesso meccanismo del lemma~\ref{lem:traiettorie} e della
        proposizione~\ref{prop:CS-invariante}.
  \item \libro{Lemma~13.2}:
        $J^*(x)\le J(\tilde z(x),x)\le\sum_v w_v(x)J^*(v)$. L'interpolazione
        dei costi nei campioni sovrastima il costo della soluzione
        interpolata, che a sua volta sovrastima l'ottimo. È la
        proposizione~\ref{prop:P}(b), e la figura~13.3 del libro è
        l'analoga della figura~\ref{fig:baricentrica}.
  \item \libro{Teorema~13.2}: per certificare la qualità
        dell'approssimazione su $\Conv(V)$ si ottimizza ``su tutte le
        possibili funzioni baricentriche'', cioè sui pesi $\lambda_v\ge 0$
        con $\sum\lambda_v=1$ e $\sum\lambda_v v=x$. Compare cioè lo stesso
        minimo sui pesi che definisce $P^j$ nel paper.
  \item \libro{Teorema~13.1}: un MPC risolto in modo \emph{subottimo} resta
        stabile se l'errore di subottimalità è minore del costo di stadio. La
        dimostrazione è ancora quella della sequenza traslata.
\end{itemize}

\paragraph{Come il paper usa questi strumenti.}
La corrispondenza è diretta se si legge ogni stato memorizzato $x^i_s$ come un
campione $v$, e la coda di traiettoria memorizzata che parte da lì,
$(u^i_s,u^i_{s+1},\dots)$, come la ``soluzione nel campione''. Allora:
\begin{itemize}
  \item l'ingresso $\bar u=\sum\lambda^{*,i}_s u^i_s$ usato nella prova di
        fattibilità (proposizione~\ref{prop:CS-invariante}) è
        l'interpolazione baricentrica~(13.18) del primo ingresso delle code
        memorizzate;
  \item la prova del Lemma~13.2 non usa l'ottimalità delle soluzioni nei
        campioni, solo la loro ammissibilità e la convessità del costo; per
        questo funziona anche con le traiettorie del paper, che sono
        ammissibili ma non ottime. Al posto del costo ottimo nel campione
        compare il costo realizzato $J^i_{s\to\infty}$, e si ottiene
        $P^j\ge\Jstar$;
  \item la scelta dei pesi che minimizza il costo interpolato (il programma
        lineare~\eqref{eq:LP}) è quella che il Teorema~13.2 usa come
        rilassamento.
\end{itemize}

\begin{table}[h]
\centering
\caption{Interpolazione baricentrica nel libro (cap.~13) e nel paper.}\label{tab:baricentrica}
\small
\begin{tabular}{p{2.6cm}p{5.6cm}p{5.9cm}}
\toprule
 & Libro, cap.~13 & Paper \\
\midrule
Scopo & ridurre il calcolo online, approssimando un MPC già progettato
      & migliorare le prestazioni, imparando insieme e costo terminali \\
Campioni & punti scelti offline (triangolazioni, griglie, Algoritmo~13.1)
         & stati visitati dal sistema in anello chiuso \\
Valori nei campioni & soluzioni \emph{ottime} $z^*(v)$ e costi ottimi
                    & traiettorie \emph{realizzate} e costi realmente pagati \\
Pesi & funzioni baricentriche fissate per regione (Def.~13.3, Lemma~13.5)
     & scelti a ogni istante da un LP che minimizza il costo \\
Che cosa si interpola & la legge di controllo
                      & il costo terminale (e implicitamente l'ingresso $\bar u$) \\
Garanzia & Teorema~13.1 (MPC subottimo)
         & Teorema~12.2 di [BBM] / Teorema~1 del paper \\
\bottomrule
\end{tabular}
\end{table}

\subsection{Programmazione multiparametrica: la forma di $P^j$ e di $\JL$}\label{sec:libro-mp}

La parte~II del libro studia i problemi di ottimo in cui un vettore di
parametri $x$ entra nei vincoli: i \emph{programmi multiparametrici}. Due
risultati si applicano al paper.

\paragraph{$P^j$ è la funzione valore di un mp-LP.}
Il problema~\eqref{eq:LP} che definisce $P^j$ ha la forma~(6.11) di [BBM],
\[
  \min_z\ c^\top z\quad\text{s.t. } Gz\le w+Sx,
\]
con variabile $z=\lambda$, parametro $x$, costo $c=\vett{c}^j$ (i costi
realizzati) e vincoli $\lambda\ge 0$, $\vett{1}^\top\lambda=1$,
$\vett{D}^j\lambda=x$ (le uguaglianze si scrivono come coppie di
disuguaglianze). Il Teorema~6.5 e il Corollario~5.1 del libro dicono che la
funzione valore di un mp-LP è \emph{convessa, continua e affine a tratti} sul
suo dominio, che qui è $\CS^j$. Si ritrova così, in forma più forte, la
proposizione~\ref{prop:P}(a), e si rende rigorosa l'immagine geometrica della
figura~\ref{fig:baricentrica}: le \emph{regioni critiche} dell'mp-LP, in cui
l'insieme dei pesi attivi non cambia, sono le proiezioni delle facce
inferiori dell'inviluppo convesso dei punti $(x^i_s,J^i_{s\to\infty})$, e su
ciascuna $P^j$ è affine.

\paragraph{$\JL$ è la funzione valore di un programma parametrico convesso.}
Il problema~\eqref{eq:lmpc-unico} ha come parametro lo stato misurato
$x^j_t$ ed è convesso congiuntamente nelle variabili di decisione e nel
parametro. La sua funzione valore $\JL$ è quindi convessa: la dimostrazione
è la stessa della proposizione~\ref{prop:P}(a) e del Teorema~6.7 del libro
(si combinano convessamente le soluzioni ottime in due stati). Il Teorema~6.7
non si applica alla lettera, perché richiede una matrice hessiana definita
positiva, mentre qui il costo è lineare nei pesi $\lambda$; per la continuità
servono i risultati generali del capitolo~5 (Teorema~5.2), oppure il rimando a
Mayne et al.\ del paper. Una conseguenza concettuale: come l'MPC del libro, anche
il LMPC di un'iterazione avrebbe in linea di principio una forma esplicita,
una legge affine a tratti dello stato (capitoli~11--12 di [BBM]). Calcolarla
non avrebbe senso, perché cambia a ogni iterazione e la sua complessità cresce
con i dati: per il LMPC la soluzione online del QP (capitolo~14 di [BBM]) è la
scelta naturale.

\subsection{Orizzonte infinito, CLQR e ottimalità}\label{sec:libro-clqr}

\paragraph{Iterazione sulla funzione valore.}
La sezione~7.4.1 del libro descrive la \emph{value function iteration}: si
parte da una stima della funzione valore su una regione, la si migliora
applicando l'equazione di Bellman~\eqref{eq:bellman} in ogni punto di una
griglia dello spazio di stato, e si interpola tra i punti della griglia. Il
LMPC si può leggere come una variante di questa idea guidata dai dati, con
tre differenze: i campioni non sono una griglia (il cui numero di punti cresce
in modo esponenziale con la dimensione dello stato) ma gli stati che il
sistema visita davvero; la stima della funzione valore non si ottiene
risolvendo l'equazione di Bellman punto per punto ma \emph{misurando} i costi
lungo le traiettorie; e l'interpolazione è convessa, il che garantisce che la
stima resti per eccesso (proposizione~\ref{prop:P}(b)) e decresca con le
iterazioni (proposizione~\ref{prop:P}(d)). È un'analogia, non
un'equivalenza: il paper non presenta il LMPC in questi termini.

\paragraph{Il costo terminale come costo della coda.}
Commentando la scelta $\Xf=\mathcal{O}^{\mathrm{LQR}}_\infty$,
$p(x)=x^\top P_\infty x$, il libro osserva che i due termini del costo
dell'MPC ``riflettono il costo a orizzonte infinito'': la somma sugli $N$
passi è il costo del tratto iniziale, in cui i vincoli possono essere attivi,
e il costo terminale è il costo esatto della \emph{coda} infinita, percorsa
dentro un insieme in cui i vincoli non sono più attivi
(equazioni~(12.24)--(12.26) di [BBM]). Nel LMPC $P^{j-1}$ ha la stessa
funzione, con due differenze: è definito su una regione molto più grande, in
cui i vincoli possono essere attivi, ed è il costo di code \emph{realmente
percorse} invece che di una coda calcolata analiticamente.

\paragraph{Orizzonte finito che restituisce l'ottimo infinito.}
Il Teorema~11.3 del libro (Sznaier e Damborg, Scokaert e Rawlings) afferma che
il problema a orizzonte finito con costo terminale $x^\top P_\infty x$ e senza
vincolo terminale coincide con il CLQR a orizzonte infinito, se lo stato
finale predetto entra ``spontaneamente'' nell'insieme in cui quel costo
terminale è esatto. Il Teorema~3 del paper ha lo stesso spirito: un orizzonte
finito restituisce l'ottimo a orizzonte infinito quando il costo terminale
descrive correttamente il costo della coda. La differenza è che nel libro la
correttezza del costo terminale è garantita da un calcolo, mentre nel LMPC è
il punto d'arrivo delle iterazioni.

\paragraph{Il riferimento dell'esempio.}
La soluzione esatta del CLQR con cui il paper confronta il LMPC
(sezione~\ref{sec:esempio}) è calcolata con l'algoritmo che il paper attribuisce
a Borrelli (2003) e che nel libro è descritto nella sezione~11.3.5: si parte
dall'insieme $\mathcal{O}^{\mathrm{LQR}}_\infty$ e si esplorano le regioni
vicine risolvendo problemi a orizzonte crescente, fino a coprire l'insieme di
stati iniziali di interesse. Il Teorema~11.4 ne descrive il risultato: una
legge di controllo continua e affine a tratti su una partizione poliedrica.

\subsection{Il doppio integratore nel libro}\label{sec:libro-esempio}

Il sistema dell'esempio del paper, il doppio integratore
$x_{k+1}=\left[\begin{smallmatrix}1&1\\0&1\end{smallmatrix}\right]x_k+
\left[\begin{smallmatrix}0\\1\end{smallmatrix}\right]u_k$, è il sistema di
riferimento del libro (equazione~(11.23) di [BBM]), con limiti numerici
diversi da quelli del paper. Tre esempi del libro aiutano a leggere il paper:
\begin{itemize}
  \item \libro{Esempio~11.1} mostra come si propagano all'indietro gli insiemi
        ammissibili $\Xc_i$, con insieme terminale controllo-invariante oppure
        senza: è l'illustrazione dei Teoremi~11.1--11.2, quindi della
        regione $K_N(\CS^{j-1})$ della sezione~\ref{sec:libro-insiemi};
  \item \libro{Esempio~11.2} calcola la legge ottima esplicita, affine a
        tratti, di un CLQR a orizzonte finito: è il tipo di soluzione che il
        paper usa come termine di confronto;
  \item \libro{Esempio~12.1} applica il receding horizon con orizzonte breve e
        \emph{senza insieme terminale}: da uno stato iniziale il sistema
        converge, da uno stato vicino il problema di ottimo diventa
        impossibile dopo pochi passi, perché le predizioni ad anello aperto
        non coincidono con l'evoluzione in anello chiuso. È esattamente il
        fallimento che un insieme terminale controllo-invariante impedisce,
        e quindi ciò che il safe set convesso garantisce al LMPC.
\end{itemize}

\subsection{Dal MIP al QP: un rilassamento che non perde nulla}\label{sec:libro-mip}

Nel LMPC non lineare del 2016 il vincolo terminale $x_{t+N|t}\in\SSet^{j-1}$
è un vincolo \emph{logico}: lo stato finale deve coincidere con il primo punto
memorizzato, \emph{oppure} con il secondo, \emph{oppure} con il terzo, e così
via. Il capitolo~16 del libro (sezione~16.4) mostra come vincoli logici di
questo tipo si traducano in disuguaglianze con variabili binarie, e il
capitolo~17 (sezione~17.4) come l'MPC che ne risulta richieda di risolvere un
MIQP. La sezione~17.4.1 ricorda che questi problemi sono NP-completi e
descrive il \emph{branch and bound}: il primo passo dell'algoritmo è
\emph{rilassare} le variabili binarie, lasciandole variare con continuità in
$[0,1]$, e risolvere il problema continuo che ne risulta; poi si ramifica
fissando le variabili una alla volta.

Il LMPC convesso del paper è esattamente quel primo rilassamento: le binarie
$\delta^i_s\in\{0,1\}$ diventano i pesi $\lambda^i_s\in[0,1]$. In generale il
rilassamento di un MIP fornisce solo un limite inferiore e una soluzione che
può non rispettare i vincoli originali; per questo il branch and bound deve
proseguire. Il punto del paper è che, \emph{per sistemi lineari}, la soluzione
del problema rilassato è già ``sicura'': lo stato finale è in $\CS^{j-1}$, che
è controllo-invariante, e tutte le garanzie restano valide. Non serve
ramificare: basta un solo QP. Il problema rilassato non è equivalente al MIP
(il suo ottimo è in generale migliore, perché ha più libertà), ma ha le stesse
proprietà, ed è questo che spiega il risparmio di tempo osservato
nell'esempio.

\subsection{Che cosa aggiunge il paper rispetto al libro}\label{sec:libro-novita}

\begin{itemize}
  \item \emph{Ingredienti terminali dai dati.} Nel libro insieme e costo
        terminali si progettano (LQR e massimo insieme invariante) con
        operazioni sui poliedri; nel paper si leggono dalle traiettorie
        registrate, senza alcuna operazione sugli insiemi.
  \item \emph{La dimensione iterativa.} Il libro tratta un singolo problema
        di controllo; non contiene l'idea di eseguire il compito più volte e
        di garantire che il costo non peggiori da un'esecuzione all'altra
        (Teoremi~2 e~3 del paper).
  \item \emph{Un uso nuovo dell'interpolazione baricentrica.} Nel libro
        serve ad approssimare una legge già nota; nel paper serve a costruire
        un costo terminale da dati non ottimi, sfruttando che la prova del
        Lemma~13.2 richiede solo ammissibilità e convessità.
  \item \emph{Dove guardare per estendere il paper.} Gli strumenti per
        rimuovere le ipotesi più forti sono nel libro: insiemi invarianti
        robusti e MPC robusto (sezioni~10.3--10.4 e capitolo~15) per
        disturbi e incertezze di modello; sistemi ibridi e MIQP
        (capitoli~16--17) per uscire dalla linearità, con il costo
        computazionale che il paper ha voluto evitare.
\end{itemize}

% ===========================================================================
\section{Ipotesi, limiti e osservazioni}\label{sec:limiti}
% ===========================================================================

\begin{itemize}
  \item \emph{Modello esatto e assenza di disturbi.} Tutta la prova di
        fattibilità ricorsiva usa la~\eqref{eq:modello-esatto}: lo stato
        realizzato coincide con quello predetto. Con incertezze di modello o
        disturbi, i punti memorizzati non sono più traiettorie esatte del
        modello di predizione e le garanzie cadono; servono versioni robuste,
        costruite con gli insiemi invarianti robusti e l'MPC robusto del libro
        \libro{Sez.~10.3--10.4 e Cap.~15}.
  \item \emph{Linearità.} È l'ipotesi che rende lecita la convessificazione
        (lemma~\ref{lem:traiettorie}). Per sistemi non lineari si torna al
        safe set discreto, oppure si usano approssimazioni locali (lavori
        successivi degli stessi autori), rinunciando a parte delle garanzie.
  \item \emph{Stesso stato iniziale.} Le iterazioni partono tutte da $x_S$. È
        ciò che rende la traiettoria precedente una candidata ammissibile
        all'istante iniziale (passo 1 del teorema~\ref{teo:1} e prima
        disuguaglianza del teorema~\ref{teo:2}).
  \item \emph{Serve una prima traiettoria ammissibile}
        (ipotesi~\ref{ip:iterazione-zero}). Il LMPC migliora, ma non crea dal
        nulla la prima soluzione.
  \item \emph{Crescita dei dati.} Il numero di variabili $\lambda$ cresce con
        il numero di stati memorizzati: il QP diventa sempre più grande. Il
        paper non discute come limitarlo.
  \item \emph{Iterazioni finite.} La teoria è scritta per traiettorie infinite;
        la versione pratica tronca le traiettorie quando lo stato ha
        raggiunto $x_F$ con tolleranza $\varepsilon$, il che introduce piccole
        approssimazioni non analizzate.
  \item \emph{Convergenza delle traiettorie.} Il teorema di ottimalità la
        assume, non la dimostra.
\end{itemize}

% ===========================================================================
\section{Riepilogo}\label{sec:riepilogo}
% ===========================================================================

\begin{table}[h]
\centering
\caption{MPC classico e LMPC a confronto.}\label{tab:confronto}
\begin{tabular}{p{3.4cm}p{5.2cm}p{5.6cm}}
\toprule
 & MPC classico & LMPC \\
\midrule
Insieme terminale & progettato (es.\ invariante LQR), spesso piccolo & $\CS^{j-1}$, inviluppo convesso degli stati visitati; cresce con le iterazioni \\
Costo terminale & progettato (es.\ $x^\top P_\infty x$) & $P^{j-1}$, interpolazione convessa dei costi realizzati; diminuisce con le iterazioni \\
Perché soddisfano la~\eqref{eq:clf} & per costruzione analitica & lemma~\ref{lem:decrescita}: linearità + convessità + costi realmente pagati \\
Problema a ogni istante & QP & QP con in più i pesi $\lambda$ \\
Garanzie & fattibilità ricorsiva, stabilità & le stesse, più miglioramento tra iterazioni e ottimalità della traiettoria limite \\
\bottomrule
\end{tabular}
\end{table}

\noindent
I passaggi logici del paper, in sintesi:
\begin{enumerate}
  \item Un compito iterativo fornisce, a ogni esecuzione, una traiettoria
        ammissibile e il costo realmente pagato da ogni suo stato.
  \item Gli stati visitati sono sicuri (proposizione~\ref{prop:SS}). Per un
        sistema lineare con vincoli convessi lo è anche ogni loro
        combinazione convessa: il safe set convesso $\CS^j$ è
        controllo-invariante (proposizione~\ref{prop:CS-invariante}).
  \item Interpolando convessamente i costi realizzati si ottiene $P^j$, una
        funzione convessa che sovrastima il costo ottimo e che decresce
        almeno del costo di stadio lungo opportune combinazioni degli ingressi
        memorizzati (lemma~\ref{lem:decrescita}).
  \item Usando $(\CS^{j-1},P^{j-1})$ come ingredienti terminali, l'MPC è un
        problema convesso (QP) e le dimostrazioni classiche danno fattibilità
        ricorsiva e stabilità (teorema~\ref{teo:1}).
  \item Il costo di ogni iterazione è limitato dal costo dell'iterazione
        precedente (teorema~\ref{teo:2}) e, se le traiettorie convergono, il
        limite è l'ottimo a orizzonte infinito (teorema~\ref{teo:3}).
  \item Rispetto al LMPC non lineare, il passaggio dall'insieme discreto
        all'inviluppo convesso trasforma un MIP in un QP, con un risparmio
        computazionale di diversi ordini di grandezza.
  \item Nel quadro del libro di Borrelli, Bemporad e Morari
        (sezione~\ref{sec:libro}): il Teorema~1 è il Teorema~12.2 di [BBM]
        con $(\Xf,p)=(\CS^{j-1},P^{j-1})$; $P^j$ è l'interpolazione
        baricentrica del capitolo~13, applicata a traiettorie realizzate
        invece che a soluzioni ottime, ed è la funzione valore di un mp-LP
        (capitolo~6); il passaggio dal MIP al QP è il rilassamento continuo
        con cui inizia il branch and bound (capitolo~17), che per sistemi
        lineari conserva tutte le garanzie.
\end{enumerate}

% ===========================================================================
\begin{thebibliography}{9}
\bibitem{rosolia2017} U.~Rosolia, F.~Borrelli, \emph{Learning Model
  Predictive Control for Iterative Tasks: A Computationally Efficient Approach
  for Linear System}, IFAC PapersOnLine 50-1 (2017), 3142--3147.
\bibitem{rosolia2016} U.~Rosolia, F.~Borrelli, \emph{Learning Model
  Predictive Control for Iterative Tasks}, arXiv:1609.01387 (2016); versione
  estesa: \emph{Learning Model Predictive Control for Iterative Tasks. A
  Data-Driven Control Framework}, IEEE Transactions on Automatic Control
  63(7) (2018).
\bibitem{mayne2000} D.Q.~Mayne, J.B.~Rawlings, C.V.~Rao, P.O.M.~Scokaert,
  \emph{Constrained model predictive control: Stability and optimality},
  Automatica 36(6) (2000), 789--814.
\bibitem{borrelli2003} F.~Borrelli, \emph{Constrained Optimal Control of
  Linear and Hybrid Systems}, Springer (2003).
\bibitem{bbm2017} F.~Borrelli, A.~Bemporad, M.~Morari, \emph{Predictive
  Control for Linear and Hybrid Systems}, Cambridge University Press (2017).
\bibitem{jones2010} C.N.~Jones, M.~Morari, \emph{Polytopic approximation of
  explicit model predictive controllers}, IEEE Transactions on Automatic
  Control 55(11) (2010), 2542--2553.
\bibitem{bristow2006} D.A.~Bristow, M.~Tharayil, A.G.~Alleyne, \emph{A survey
  of iterative learning control}, IEEE Control Systems 26(3) (2006), 96--114.
\bibitem{boyd2004} S.~Boyd, L.~Vandenberghe, \emph{Convex Optimization},
  Cambridge University Press (2004).
\end{thebibliography}

\end{document}
